\documentclass[11pt,oneside]{scrartcl}

% -----------------------------------------------------------------------------
% Idioma, tipografía y composición
% -----------------------------------------------------------------------------
\usepackage[utf8]{inputenc}
\usepackage[T1]{fontenc}
\usepackage[spanish,es-nodecimaldot]{babel}
\usepackage{amsmath,amsthm,mathtools,bm}
\usepackage{newtxtext}
\usepackage{newtxmath}
\usepackage{microtype}
\usepackage[a4paper,top=2.55cm,bottom=2.65cm,left=2.55cm,right=2.55cm,
            headheight=17pt,headsep=0.65cm,footskip=1.15cm]{geometry}
\usepackage{setspace}
\setstretch{1.055}
\setlength{\parindent}{1.25em}
\setlength{\parskip}{0.20em}
\usepackage{ragged2e}
\usepackage{scrlayer-scrpage}
\clearpairofpagestyles
\ihead{\small\sffamily De la conjugación armónica a una torre polar de medias}
\ohead{\small\sffamily Robert Duan Montoya Cardona}
\cfoot{\small\sffamily\thepage}
\setkomafont{pageheadfoot}{\normalfont}
\setkomafont{section}{\sffamily\bfseries\color{Navy}}
\setkomafont{subsection}{\sffamily\bfseries\color{DeepBlue}}
\setkomafont{subsubsection}{\sffamily\bfseries\color{Slate}}
\setkomafont{paragraph}{\sffamily\bfseries\color{Slate}}

% -----------------------------------------------------------------------------
% Matemáticas
% -----------------------------------------------------------------------------
\usepackage{thmtools}
\usepackage{array,booktabs,longtable,tabularx,multirow}
\usepackage{enumitem}
\usepackage{siunitx}
\sisetup{output-decimal-marker={,}}

% -----------------------------------------------------------------------------
% Gráficos y color
% -----------------------------------------------------------------------------
\usepackage{xcolor}
\definecolor{Navy}{HTML}{132A46}
\definecolor{DeepBlue}{HTML}{1D4E89}
\definecolor{Sky}{HTML}{DCECF7}
\definecolor{Ice}{HTML}{F3F8FC}
\definecolor{Gold}{HTML}{B88322}
\definecolor{Warm}{HTML}{F6E9D1}
\definecolor{Brick}{HTML}{9B3A32}
\definecolor{Green}{HTML}{2E6B57}
\definecolor{Mint}{HTML}{DDEFE8}
\definecolor{Slate}{HTML}{485769}
\definecolor{SoftGray}{HTML}{EFF2F5}
\definecolor{DarkGray}{HTML}{30363D}
\usepackage{tikz}
\usetikzlibrary{arrows.meta,calc,positioning,fit,backgrounds,decorations.pathreplacing,
                intersections,patterns,shadows.blur,matrix,3d,angles,quotes,babel}
\usepackage{tikz-cd}
\usepackage{pgfplots}
\pgfplotsset{compat=1.18}
\usepackage{caption}
\usepackage{subcaption}
\usepackage[section]{placeins}
\usepackage{needspace}
\captionsetup{font=small,labelfont={bf,sf,color=Navy},textfont=it,
              justification=justified,singlelinecheck=false}

% -----------------------------------------------------------------------------
% Cajas, enlaces y referencias internas
% -----------------------------------------------------------------------------
\usepackage[most]{tcolorbox}
\usepackage{csquotes}
\usepackage{hyperref}
\hypersetup{
  colorlinks=true,
  linkcolor=DeepBlue,
  citecolor=Green,
  urlcolor=Brick,
  pdftitle={De la conjugación armónica a una torre polar de medias},
  pdfauthor={Robert Duan Montoya Cardona},
  pdfsubject={Geometría proyectiva, apolaridad, medias simétricas y dinámica racional},
  pdfkeywords={apolaridad, curva racional normal, medias simétricas, Newton, mapa polar, formas binarias}
}
\usepackage[capitalise,noabbrev,nameinlink]{cleveref}
\crefname{theorem}{teorema}{teoremas}
\crefname{proposition}{proposición}{proposiciones}
\crefname{lemma}{lema}{lemas}
\crefname{corollary}{corolario}{corolarios}
\crefname{definition}{definición}{definiciones}
\crefname{remark}{observación}{observaciones}
\crefname{example}{ejemplo}{ejemplos}
\crefname{figure}{figura}{figuras}
\crefname{section}{sección}{secciones}
\crefname{equation}{ecuación}{ecuaciones}

% -----------------------------------------------------------------------------
% Entornos matemáticos
% -----------------------------------------------------------------------------
\declaretheoremstyle[
  headfont=\sffamily\bfseries\color{Navy},
  bodyfont=\itshape,
  spaceabove=8pt,spacebelow=8pt,
  headpunct={.},postheadspace=0.55em
]{plainblue}
\declaretheoremstyle[
  headfont=\sffamily\bfseries\color{DeepBlue},
  bodyfont=\normalfont,
  spaceabove=8pt,spacebelow=8pt,
  headpunct={.},postheadspace=0.55em
]{definitionblue}
\declaretheoremstyle[
  headfont=\sffamily\bfseries\color{Slate},
  bodyfont=\normalfont,
  spaceabove=7pt,spacebelow=7pt,
  headpunct={.},postheadspace=0.55em
]{remarkgray}

\theoremstyle{plainblue}
\newtheorem{theorem}{Teorema}[section]
\newtheorem{proposition}[theorem]{Proposición}
\newtheorem{lemma}[theorem]{Lema}
\newtheorem{corollary}[theorem]{Corolario}
\theoremstyle{definitionblue}
\newtheorem{definition}[theorem]{Definición}
\newtheorem{example}[theorem]{Ejemplo}
\theoremstyle{remarkgray}
\newtheorem{remark}[theorem]{Observación}
\newtheorem{warning}[theorem]{Advertencia}

% Identificadores PDF unívocos para los entornos que comparten contador.
\makeatletter
\def\theHtheorem{\thesection.\arabic{theorem}}
\def\theHproposition{\theHtheorem}
\def\theHlemma{\theHtheorem}
\def\theHcorollary{\theHtheorem}
\def\theHdefinition{\theHtheorem}
\def\theHexample{\theHtheorem}
\def\theHremark{\theHtheorem}
\def\theHwarning{\theHtheorem}
\makeatother

\newtcolorbox{keybox}[1][]{
  enhanced,breakable,colback=Ice,colframe=DeepBlue,boxrule=0.8pt,
  arc=2mm,left=4mm,right=4mm,top=2.5mm,bottom=2.5mm,
  borderline west={2.5pt}{0pt}{DeepBlue},fonttitle=\sffamily\bfseries,
  title=#1
}
\newtcolorbox{researchbox}[1][]{
  enhanced,breakable,colback=Warm,colframe=Gold,boxrule=0.8pt,
  arc=2mm,left=4mm,right=4mm,top=2.5mm,bottom=2.5mm,
  borderline west={2.5pt}{0pt}{Gold},fonttitle=\sffamily\bfseries,
  title=#1
}
\newtcolorbox{proofmap}[1][]{
  enhanced,breakable,colback=SoftGray,colframe=Slate,boxrule=0.5pt,
  arc=1.5mm,left=4mm,right=4mm,top=2mm,bottom=2mm,
  fonttitle=\sffamily\bfseries,title=#1
}

% -----------------------------------------------------------------------------
% Operadores y abreviaturas
% -----------------------------------------------------------------------------
\newcommand{\RR}{\mathbb{R}}
\newcommand{\CC}{\mathbb{C}}
\newcommand{\QQ}{\mathbb{Q}}
\newcommand{\KK}{\mathbb{K}}
\newcommand{\PP}{\mathbb{P}}
\newcommand{\Aff}{\mathbb{A}}
\newcommand{\ee}{\mathrm{e}}
\newcommand{\dd}{\mathrm{d}}
\newcommand{\id}{\mathrm{id}}
\newcommand{\Sym}{\operatorname{Sym}}
\newcommand{\PGL}{\operatorname{PGL}}
\newcommand{\GL}{\operatorname{GL}}
\newcommand{\SL}{\operatorname{SL}}
\newcommand{\Osc}{\operatorname{Osc}}
\newcommand{\Span}{\operatorname{span}}
\newcommand{\Fix}{\operatorname{Fix}}
\newcommand{\Crit}{\operatorname{Crit}}
\newcommand{\rad}{\operatorname{rad}}
\newcommand{\Res}{\operatorname{Res}}
\newcommand{\Hess}{\operatorname{Hess}}
\newcommand{\AM}{\operatorname{AM}}
\newcommand{\GM}{\operatorname{GM}}
\newcommand{\HM}{\operatorname{HM}}
\newcommand{\CR}{\operatorname{cr}}
\newcommand{\cN}{\mathfrak{N}}
\newcommand{\cP}{\mathcal{P}}
\newcommand{\cD}{\mathcal{D}}
\newcommand{\cM}{\mathcal{M}}
\newcommand{\cE}{\mathcal{E}}
\newcommand{\cC}{\mathcal{C}}
\newcommand{\cF}{\mathcal{F}}
\newcommand{\cH}{\mathcal{H}}
\newcommand{\cO}{\mathcal{O}}
\newcommand{\brak}[1]{\left[#1\right]}
\newcommand{\paren}[1]{\left(#1\right)}
\newcommand{\set}[1]{\left\{#1\right\}}
\newcommand{\abs}[1]{\left|#1\right|}
\newcommand{\norm}[1]{\left\lVert#1\right\rVert}
\newcommand{\xto}[1]{\xrightarrow{#1}}
\newcommand{\defeq}{\vcentcolon=}

% -----------------------------------------------------------------------------
% Datos del documento
% -----------------------------------------------------------------------------
\title{De la conjugación armónica a una torre polar de medias}
\subtitle{Apolaridad, derivadas normalizadas, formas binarias y dinámica racional de divisores de raíces}
\author{Robert Duan Montoya Cardona}
\date{Julio de 2026}

\begin{document}
\hypersetup{pageanchor=false}

% -----------------------------------------------------------------------------
% Portada
% -----------------------------------------------------------------------------
\begin{titlepage}
\thispagestyle{empty}
\begin{tikzpicture}[remember picture,overlay]
  \fill[Navy] (current page.north west) rectangle ([yshift=-9.30cm]current page.north east);
  \fill[Ice] ([yshift=-9.30cm]current page.north west) rectangle (current page.south east);
  \draw[Gold,line width=1.2pt] ([xshift=1.7cm,yshift=-9.15cm]current page.north west)
       -- ([xshift=-1.7cm,yshift=-9.15cm]current page.north east);

  % Curva racional normal esquemática
  \begin{scope}[shift={(13.0,-2.05)},scale=0.72]
    \draw[Sky,line width=9pt,opacity=.24] plot[smooth] coordinates
      {(-4,-.6) (-3.3,.15) (-2.5,.65) (-1.5,.75) (-.4,.25) (.8,-.6) (2,-1.0) (3.1,-.55) (4,.35)};
    \draw[white,line width=1.5pt] plot[smooth] coordinates
      {(-4,-.6) (-3.3,.15) (-2.5,.65) (-1.5,.75) (-.4,.25) (.8,-.6) (2,-1.0) (3.1,-.55) (4,.35)};
    \foreach \x/\y in {-3.2/.23,-1.45/.73,.65/-.50,2.1/-.97,3.25/-.45}{
      \fill[Gold] (\x,\y) circle (2.2pt);
    }
    \draw[white!80,dashed] (-3.2,.23)--(3.25,-.45);
  \end{scope}

  % Cadena de medias
  \begin{scope}[shift={(7.10,-17.8)}]
    \foreach \i/\lab in {0/{\AM},1/{\cN_2},2/{\cdots},3/{\cN_{n-1}},4/{\HM}}{
      \node[circle,draw=DeepBlue,fill=white,minimum size=12mm,
            font=\small\bfseries,blur shadow] (m\i) at (1.75*\i,0) {$\lab$};
    }
    \foreach \i in {0,...,3}{\draw[-{Latex[length=2mm]},DeepBlue,thick] (m\i)--(m\the\numexpr\i+1\relax);}
    \node[font=\small\sffamily,text=Slate] at (3.5,-1.15)
      {$\displaystyle \prod_{k=1}^{n}\cN_k=\GM^n$};
  \end{scope}
\end{tikzpicture}

\vspace*{1.0cm}
{\color{white}\sffamily
\begin{center}
  {\fontsize{27}{31}\selectfont\bfseries De la conjugación armónica\\[2mm]
  a una torre polar de medias\par}
  \vspace{0.65cm}
  {\Large Apolaridad, derivadas normalizadas, formas binarias\\
  y dinámica racional de divisores de raíces\par}
\end{center}}

\vspace{4.45cm}
\begin{center}
{\Large\sffamily\bfseries\color{Navy} Robert Duan Montoya Cardona}\par
\vspace{0.35cm}
{\large\sffamily\color{Slate} Manuscrito de investigación autocontenido}\par
\vspace{0.15cm}
{\sffamily\color{Slate} Geometría proyectiva \enspace$\cdot$\enspace Álgebra clásica \enspace$\cdot$\enspace Desigualdades simétricas \enspace$\cdot$\enspace Dinámica compleja}\par
\vfill
\begin{minipage}{0.82\textwidth}
\centering\small\color{DarkGray}
Esta versión reconstruye desde sus fundamentos la relación entre la parábola, la
conjugación armónica, la curva racional normal y la apolaridad, y la prolonga hacia
una torre de mapas polares asociada a las derivadas normalizadas de un polinomio.
El texto distingue cuidadosamente los ingredientes clásicos de las contribuciones
estructurales desarrolladas aquí.
\end{minipage}\par
\vspace{0.75cm}
{\sffamily\color{Slate} Julio de 2026\par}
\end{center}
\end{titlepage}
\hypersetup{pageanchor=true}

\pagenumbering{roman}

\begin{abstract}
Sea
\[
P(t)=\prod_{i=1}^{n}(t-r_i)
=t^n-e_1t^{n-1}+e_2t^{n-2}-\cdots+(-1)^ne_n
\]
un polinomio mónico con raíces positivas. El propósito de este trabajo es explicar,
desde los fundamentos de la recta proyectiva hasta la geometría de formas binarias,
por qué las razones consecutivas
\[
\cN_k(r_1,\dots,r_n)
=\frac{e_k/\binom nk}{e_{k-1}/\binom n{k-1}}
=\frac{k}{n-k+1}\frac{e_k}{e_{k-1}},
\qquad 1\le k\le n,
\]
forman una cadena natural de medias que comienza en la media aritmética y termina
en la media armónica. Se demuestra que
\[
\AM=\cN_1\ge \cN_2\ge\cdots\ge\cN_n=\HM,
\qquad
\prod_{k=1}^{n}\cN_k=\GM^n,
\]
y que cada \(\cN_k\) es una media simétrica, homogénea, interna y estrictamente
creciente en cada variable positiva. Estas cantidades se realizan como cocientes de
coordenadas consecutivas del punto apolar \(K_P\), obtenido como intersección de los
hiperplanos osculadores de la curva racional normal en las raíces.

La segunda estructura central es la torre de polinomios mónicos
\[
Q_k(t)=\frac{k!}{n!}P^{(n-k)}(t),\qquad Q_k'=kQ_{k-1},
\]
y sus mapas polares
\[
\Phi_k(t)=t-k\frac{Q_k(t)}{Q_k'(t)}
=t-\frac{Q_k(t)}{Q_{k-1}(t)}.
\]
Se prueba que \(\cN_k=\Phi_k(0)\). Si \(P\) tiene raíces reales positivas, las
raíces de \(Q_k\) también son positivas y
\(\cN_k\) es exactamente su media armónica, mientras que todas las capas de la
torre conservan la misma media aritmética. En el nivel superior,
\(\Phi_n(t)=t-nP(t)/P'(t)\) es el mapa racional asociado al diferencial exterior de
la forma binaria que codifica las raíces. Este mapa, ya presente en la literatura de
dinámica algebraica, se caracteriza como el único mapa racional de grado \(n-1\)
que fija las raíces simples con multiplicador \(1-n\). Su divisor crítico es el
hessiano de la forma binaria. Para \(n=2\), el mapa se reduce exactamente a la
involución armónica realizada por una parábola.

El resultado es una síntesis rigurosa entre geometría proyectiva, apolaridad,
desigualdades de Newton, teoría clásica de invariantes y dinámica racional. Se
incluyen ejemplos, diagramas, extensiones a raíces múltiples y una agenda de
problemas abiertos formulada sin atribuir novedad a los componentes clásicos.
\end{abstract}

\begin{keybox}[Lectura de alcance y prioridad]
La curva racional normal, la apolaridad de formas binarias, las desigualdades de
Newton--Maclaurin y el mapa racional obtenido del diferencial de una forma binaria
son ingredientes clásicos. En particular, Doyle y McMullen observaron que el mapa
\(t-nP/P'\) es el único de grado \(n-1\) que fija las raíces con multiplicador
\(1-n\). La contribución estructural desarrollada en este manuscrito consiste en
organizar esos ingredientes en un único marco: la cadena proyectiva de cocientes
normalizados \(\cN_k\), su realización mediante banderas osculadoras opuestas, la
torre derivada \(Q_k\), la identidad \(\cN_k=\Phi_k(0)\), la interpretación de cada
capa como media armónica de raíces derivadas y el producto telescópico que recupera
la media geométrica. No se formula una afirmación absoluta de prioridad histórica
sin una búsqueda bibliográfica especializada exhaustiva.
\end{keybox}

\tableofcontents
\vspace{0.5cm}
\begin{center}
\begin{tikzpicture}[node distance=8mm and 9mm,>=Latex,every node/.style={align=center}]
\node[draw=Navy,fill=Ice,rounded corners,minimum width=47mm,minimum height=12mm,font=\sffamily\bfseries] (a) {Parábola y razón doble};
\node[draw=Navy,fill=Ice,rounded corners,minimum width=47mm,minimum height=12mm,font=\sffamily\bfseries,right=of a] (b) {Curva racional normal};
\node[draw=Navy,fill=Ice,rounded corners,minimum width=47mm,minimum height=12mm,font=\sffamily\bfseries,right=of b] (c) {Punto apolar $K_P$};
\node[draw=Brick,fill=white,rounded corners,minimum width=47mm,minimum height=12mm,font=\sffamily\bfseries,below=of a] (f) {Dinámica, hessiano\\y módulos};
\node[draw=Green,fill=Mint,rounded corners,minimum width=47mm,minimum height=12mm,font=\sffamily\bfseries,below=of b] (e) {Torre $Q_k$\\y mapas $\Phi_k$};
\node[draw=Gold,fill=Warm,rounded corners,minimum width=47mm,minimum height=12mm,font=\sffamily\bfseries,below=of c] (d) {Cadena\\$\cN_1,\dots,\cN_n$};
\draw[->,thick,DeepBlue] (a)--(b);
\draw[->,thick,DeepBlue] (b)--(c);
\draw[->,thick,Gold] (c)--(d);
\draw[->,thick,Green] (d)--(e);
\draw[->,thick,Brick] (e)--(f);
\draw[->,thick,Slate,dashed] (a)--(f);
\end{tikzpicture}
\end{center}

\clearpage
\pagenumbering{arabic}

% =============================================================================
\section{Qué problema se está estudiando realmente}
% =============================================================================

\subsection{Tres medias conocidas y una pregunta menos conocida}

Para números positivos \(r,s\), las tres medias pitagóricas son
\[
\AM(r,s)=\frac{r+s}{2},\qquad
\GM(r,s)=\sqrt{rs},\qquad
\HM(r,s)=\frac{2rs}{r+s}.
\]
Su desigualdad clásica
\[
\HM(r,s)\le \GM(r,s)\le \AM(r,s)
\]
es elemental. Sin embargo, una fórmula puede ser fácil de calcular y, al mismo
tiempo, esconder una estructura difícil de ver. La pregunta de este manuscrito no es
cómo calcular esas medias, sino por qué pueden aparecer como coordenadas de una
misma construcción proyectiva.

Consideremos la parábola
\[
y=a(x-r)(x-s),\qquad a\ne 0.
\]
Su vértice tiene abscisa \((r+s)/2\), mientras que una secante especial recupera
\(2rs/(r+s)\). Esta coincidencia sugiere que las medias no están entrando como
operaciones numéricas aisladas, sino como imágenes de puntos distinguidos bajo una
transformación proyectiva.

Cuando se sustituyen dos raíces por \(n\) raíces, el plano ya no contiene por sí solo
la geometría adecuada. El objeto natural pasa a ser la curva racional normal
\[
\nu_n:\PP^1\longrightarrow\PP^n,
\qquad
[u:v]\longmapsto[u^n:u^{n-1}v:\cdots:uv^{n-1}:v^n].
\]
Las raíces determinan puntos de esta curva, sus hiperplanos osculadores se cortan en
un punto \(K_P\), y las coordenadas normalizadas de ese punto son precisamente los
polinomios simétricos elementales de las raíces.

\subsection{La nueva dirección: no sólo dos medias, sino una cadena completa}

Sea
\[
e_k=e_k(r_1,\dots,r_n)
=\sum_{1\le i_1<\cdots<i_k\le n}r_{i_1}\cdots r_{i_k},
\qquad e_0=1.
\]
Los números normalizados
\[
a_k=\frac{e_k}{\binom nk}
\]
son las medias elementales sin extraer la raíz \(k\)-ésima. En lugar de estudiar
únicamente \(a_k^{1/k}\), como en las desigualdades de Maclaurin, estudiaremos sus
cocientes consecutivos:
\begin{equation}\label{eq:def-Nk-intro}
\cN_k=\frac{a_k}{a_{k-1}}
=\frac{k}{n-k+1}\frac{e_k}{e_{k-1}}.
\end{equation}
Los extremos son
\[
\cN_1=\frac{e_1}{n}=\AM(r_1,\dots,r_n),
\qquad
\cN_n=\frac{ne_n}{e_{n-1}}=\HM(r_1,\dots,r_n).
\]
La pregunta cambia entonces de naturaleza:

\begin{researchbox}[Pregunta guía]
¿Existe una geometría que produzca simultáneamente todos los cocientes
\(\cN_1,\dots,\cN_n\), explique por qué están ordenados, muestre por qué su
producto es \(\GM^n\) y conecte esa cadena con mapas racionales definidos por las
raíces y por las derivadas sucesivas de \(P\)?
\end{researchbox}

La respuesta será afirmativa. El punto apolar \(K_P\) contiene todos los \(a_k\) como
coordenadas consecutivas. Las banderas osculadoras en \(0\) y \(\infty\) seleccionan
cada par \([a_k:a_{k-1}]\). Al mismo tiempo, las derivadas normalizadas de \(P\)
forman una torre \(Q_k\), y el mapa polar asociado a \(Q_k\), evaluado en el origen,
produce exactamente \(\cN_k\).

\subsection{Mapa conceptual de los resultados principales}

\begin{figure}[ht]
\centering
\begin{tikzpicture}[node distance=12mm and 12mm,>=Latex,scale=0.85,transform shape,
 every node/.style={align=center}]
\node[draw=Navy,fill=Ice,rounded corners,minimum width=39mm,minimum height=13mm] (roots)
  {raíces $r_1,\dots,r_n$\\[-1mm] divisor en $\PP^1$};
\node[draw=Navy,fill=Ice,rounded corners,minimum width=42mm,minimum height=13mm,right=of roots] (poly)
  {$P(t)=\prod(t-r_i)$\\ coeficientes $e_k$};
\node[draw=DeepBlue,fill=Sky,rounded corners,minimum width=43mm,minimum height=13mm,right=of poly] (kp)
  {punto apolar $K_P$\\ intersección osculadora};
\node[draw=Gold,fill=Warm,rounded corners,minimum width=48mm,minimum height=13mm,below=of kp] (chain)
  {$[a_k:a_{k-1}]$\\ cadena $\cN_k=a_k/a_{k-1}$};
\node[draw=Green,fill=Mint,rounded corners,minimum width=48mm,minimum height=13mm,left=of chain] (tower)
  {$Q_k=\frac{k!}{n!}P^{(n-k)}$\\ $Q_k'=kQ_{k-1}$};
\node[draw=Brick,fill=white,rounded corners,minimum width=48mm,minimum height=13mm,left=of tower] (map)
  {$\Phi_k(t)=t-\frac{Q_k}{Q_{k-1}}$\\ $\Phi_k(0)=\cN_k$};
\draw[->,thick,DeepBlue] (roots)--(poly);
\draw[->,thick,DeepBlue] (poly)--(kp);
\draw[->,thick,Gold] (kp)--(chain);
\draw[->,thick,Green] (poly)--(tower);
\draw[->,thick,Green] (tower)--(map);
\draw[->,thick,Brick] (map)--node[below,sloped,font=\small] {misma cantidad}(chain);
\draw[->,thick,Slate,dashed,bend right=28] (roots) to node[above,font=\small] {forma binaria}(map);
\end{tikzpicture}
\caption{Las dos rutas que convergen en la cadena de medias: la ruta apolar y la ruta dinámica.}
\label{fig:mapa-principal}
\end{figure}

% =============================================================================
\section{Contexto histórico y conceptual}
% =============================================================================

\subsection{División armónica y geometría proyectiva}

La geometría proyectiva nació de la necesidad de formular propiedades que no
cambian al proyectar una figura desde un punto de vista a otro. Distancias y ángulos
no sobreviven en general, pero la incidencia, la alineación y la razón doble sí lo
hacen. La división armónica, caracterizada por una razón doble igual a \(-1\), fue
uno de los mecanismos clásicos para construir conjugaciones sobre una recta sin
recurrir a una métrica euclidiana \cite{Coxeter,Berger}.

En lenguaje moderno, dos puntos distintos \(r,s\in\PP^1\) determinan una única
involución proyectiva no trivial que los fija. En una carta afín, esa involución es
\[
h_{r,s}(o)=\frac{2rs-(r+s)o}{r+s-2o}.
\]
La media aritmética y la media armónica aparecen simplemente como
\[
h_{r,s}(\infty)=\frac{r+s}{2},
\qquad
h_{r,s}(0)=\frac{2rs}{r+s}.
\]
La parábola proporciona una realización visual de esa involución.

\subsection{Newton, Maclaurin y la geometría oculta de los coeficientes}

Si \(r_1,\dots,r_n>0\), las desigualdades de Newton afirman que la sucesión
normalizada
\[
a_k=\frac{e_k}{\binom nk}
\]
es log-cóncava:
\[
a_k^2\ge a_{k-1}a_{k+1}.
\]
Las desigualdades de Maclaurin se expresan usualmente mediante las cantidades
\(a_k^{1/k}\) \cite{HardyLittlewoodPolya,MarshallOlkinArnold}. Sin embargo, la log-concavidad contiene otra lectura inmediata:
\[
\frac{a_1}{a_0}\ge \frac{a_2}{a_1}\ge\cdots\ge\frac{a_n}{a_{n-1}}.
\]
Es exactamente la cadena \(\cN_1\ge\cdots\ge\cN_n\). Así, el paso de Maclaurin a
los cocientes consecutivos no añade una desigualdad nueva, pero revela una
organización distinta: una escalera multiplicativa cuyas piezas tienen dimensión de
longitud y cuyos extremos son \(\AM\) y \(\HM\).

\subsection{Apolaridad, formas binarias y la curva racional normal}

La apolaridad surgió en la teoría clásica de invariantes como un emparejamiento entre
formas binarias. En la formulación geométrica contemporánea, la aplicación de
Veronese
\[
\PP^1\hookrightarrow\PP^n
\]
identifica una potencia \(L^n\) de una forma lineal con un punto de una curva racional
normal. Los hiperplanos osculadores de esa curva pueden describirse mediante la
anulación apolar \cite{Dolgachev,IarrobinoKanev,Olver}. El polinomio cuyas raíces son \(r_i\) define un hiperplano secante;
su polo apolar es un punto \(K_P\) que también es la intersección de los hiperplanos
osculadores en las raíces. La geometría de configuraciones osculadoras de curvas
racionales normales continúa siendo objeto de investigación moderna
\cite{CaminataCarliniSchaffler}.

\subsection{El diferencial exterior y los mapas racionales de raíces}

Una forma binaria homogénea \(F(X,Y)\) determina la $1$-forma exacta
\(\dd F=F_X\dd X+F_Y\dd Y\). En dimensión dos, una forma lineal puede identificarse,
mediante el área orientada, con un vector. De aquí nace el mapa homogéneo
\[
[-F_Y:F_X]:\PP^1\longrightarrow\PP^1.
\]
Para \(F(X,Y)=Y^nP(X/Y)\), su expresión afín es
\[
\Phi_P(t)=t-n\frac{P(t)}{P'(t)}.
\]
Este mapa fue señalado explícitamente por Doyle y McMullen en el contexto de mapas
racionales con simetría \cite{DoyleMcMullen}: es el único mapa de grado \(n-1\) que fija las raíces de
\(P\) y tiene derivada \(1-n\) en cada una de ellas. En el presente trabajo no se
reclama ese hecho como nuevo; se utiliza como el nivel superior de una torre que
conecta todas las medias \(\cN_k\).

\begin{figure}[ht]
\centering
\begin{tikzpicture}[x=1.2cm,y=1cm,>=Latex]
\draw[Slate,line width=1.1pt] (0,0)--(12,0);
\foreach \x/\year/\lab in {
0.5/{s. XIX}/{razón doble\\y división armónica},
3.4/{s. XIX}/{invariantes\\y apolaridad},
6.2/{Newton--Maclaurin}/{log-concavidad\\de $e_k/\binom nk$},
9.0/{1989}/{Doyle--McMullen\\mapa $t-nP/P'$},
11.6/{presente marco}/{cadena $\cN_k$\\y torre derivada}}
{
  \fill[DeepBlue] (\x,0) circle (2.4pt);
  \draw[DeepBlue] (\x,0)--(\x,0.35);
  \node[above,font=\small\sffamily\bfseries,text=Navy] at (\x,0.35) {\year};
  \node[below=3mm,align=center,font=\scriptsize,text=Slate] at (\x,0) {\lab};
}
\end{tikzpicture}
\caption{Línea conceptual, no cronología exhaustiva: el manuscrito reúne tradiciones matemáticas distintas en una sola construcción.}
\label{fig:timeline}
\end{figure}

% =============================================================================
\section{La recta proyectiva, la razón doble y la involución armónica}
% =============================================================================

\subsection{La recta proyectiva desde cero}

Sea \(\KK\) un cuerpo. La recta proyectiva se define como
\[
\PP^1(\KK)=\bigl(\KK^2\setminus\{(0,0)\}\bigr)/\sim,
\]
donde \((u,v)\sim(\lambda u,\lambda v)\) para todo \(\lambda\in\KK^\times\). La
clase de \((u,v)\) se escribe \([u:v]\). Cuando \(v\ne0\), podemos dividir por
\(v\) y escribir \([u:v]=[u/v:1]\). Así aparece la carta afín \(t=u/v\). El punto
que falta es \([1:0]\), denotado por \(\infty\).

Una transformación proyectiva de la recta es la acción de una matriz invertible
\[
\begin{pmatrix}a&b\\c&d\end{pmatrix}
\]
sobre coordenadas homogéneas. En la carta afín se convierte en una transformación de
Möbius
\[
t\longmapsto\frac{at+b}{ct+d}.
\]
Dos matrices proporcionales inducen la misma transformación, por lo que el grupo
natural es \(\PGL_2(\KK)\).

\subsection{Razón doble}

\begin{definition}[Razón doble homogénea]
Sean \(A=[a_0:a_1]\), \(B=[b_0:b_1]\), \(C=[c_0:c_1]\) y
\(D=[d_0:d_1]\) cuatro puntos tales que los denominadores siguientes no se anulan.
Definimos
\[
(A,B;C,D)=
\frac{\det(A,C)\det(B,D)}{\det(A,D)\det(B,C)},
\]
donde \(\det(A,C)=a_0c_1-a_1c_0\).
\end{definition}

La fórmula es homogénea: si se cambia un representante por un múltiplo escalar, ese
factor aparece una vez arriba y una vez abajo, y se cancela.

\begin{proposition}[Invariancia proyectiva]
La razón doble es invariante bajo la acción de \(\PGL_2(\KK)\).
\end{proposition}

\begin{proof}
Si \(M\in\GL_2(\KK)\), entonces
\[
\det(Mu,Mv)=\det(M)\det(u,v).
\]
Cada uno de los cuatro determinantes de la razón doble recibe el mismo factor
\(\det(M)\). Dos factores aparecen en el numerador y dos en el denominador, por lo
que se cancelan.
\end{proof}

\subsection{Conjugados armónicos}

\begin{definition}
Dados dos puntos distintos \(R,S\in\PP^1\), decimos que \(O\) y \(H\) son
conjugados armónicos respecto de \(R,S\) si
\[
(R,S;O,H)=-1.
\]
\end{definition}

\begin{theorem}[Fórmula de la involución armónica]\label{thm:harmonic-map}
Sean \(r,s\in\KK\), \(r\ne s\), y supongamos \(\operatorname{char}(\KK)\ne2\).
Para cada \(o\in\PP^1(\KK)\) existe un único \(h_{r,s}(o)\) tal que
\((r,s;o,h_{r,s}(o))=-1\). En coordenada afín,
\begin{equation}\label{eq:harmonic-explicit}
h_{r,s}(o)=\frac{2rs-(r+s)o}{r+s-2o}.
\end{equation}
Si
\[
A=\frac{r+s}{2},\qquad D=\frac{r-s}{2},
\]
entonces
\begin{equation}\label{eq:harmonic-centered}
h_{r,s}(o)=A+\frac{D^2}{o-A}.
\end{equation}
Además, \(h_{r,s}\) es una involución proyectiva no trivial, fija \(r\) y \(s\), y
satisface
\[
h_{r,s}(\infty)=A,
\qquad
h_{r,s}(0)=\frac{2rs}{r+s}
\]
cuando \(r+s\ne0\).
\end{theorem}

\begin{proof}
En coordenadas afines, la condición armónica es
\[
\frac{r-o}{s-o}\frac{r-h}{s-h}=-1.
\]
Al despejar \(h\),
\[
(r-o)(s-h)=-(s-o)(r-h),
\]
y, después de reagrupar,
\[
h(r+s-2o)=2rs-(r+s)o.
\]
Esto produce \cref{eq:harmonic-explicit}. Para obtener la forma centrada, sustituimos
\(r=A+D\), \(s=A-D\), de modo que \(rs=A^2-D^2\). Entonces
\[
\frac{2(A^2-D^2)-2Ao}{2A-2o}
=A+\frac{D^2}{o-A}.
\]
La involutividad se ve directamente:
\[
h_{r,s}(o)-A=\frac{D^2}{o-A}
\quad\Longrightarrow\quad
h_{r,s}(h_{r,s}(o))-A
=\frac{D^2}{D^2/(o-A)}=o-A.
\]
Las demás identidades se obtienen por sustitución o por continuidad proyectiva.
\end{proof}

\begin{figure}[ht]
\centering
\begin{tikzpicture}[scale=1.03,>=Latex]
\draw[-{Latex[length=2.5mm]},Slate] (-5.3,0)--(5.4,0) node[right] {$\PP^1$};
\coordinate (R) at (-3.4,0);
\coordinate (S) at (3.2,0);
\coordinate (O) at (-1.0,0);
\coordinate (H) at (1.6,0);
\foreach \p/\lab/\col in {R/$r$/Brick,S/$s$/Brick,O/$o$/DeepBlue,H/$h_{r,s}(o)$/Gold}{
  \fill[\col] (\p) circle (2.6pt);
  \node[above=2mm,text=\col,font=\small\bfseries] at (\p) {\lab};
}
\draw[decorate,decoration={brace,amplitude=5pt},Brick] (R)+(0,-.35)--(S)+(0,-.35)
 node[midway,below=5pt,font=\small] {par fijo};
\draw[decorate,decoration={brace,amplitude=5pt,mirror},DeepBlue] (O)+(0,.42)--(H)+(0,.42)
 node[midway,above=6pt,font=\small] {par armónico};
\node[draw=Navy,fill=Ice,rounded corners,align=center,font=\small] at (0,-1.45)
 {$\displaystyle (r,s;o,h)=-1$\\la razón doble, no la distancia, es el invariante};
\end{tikzpicture}
\caption{La condición armónica vive en la recta proyectiva y no depende de una métrica.}
\label{fig:harmonic-line}
\end{figure}

% =============================================================================
\section{La parábola como máquina geométrica de la involución}
% =============================================================================

\subsection{Construcción correcta}

Sea
\[
\mathcal P:\quad y=a(x-r)(x-s),\qquad a\in\RR^\times,\quad r\ne s.
\]
Escribamos
\[
A=\frac{r+s}{2},\qquad D=\frac{r-s}{2}.
\]
Entonces
\[
\mathcal P:\quad y=a\bigl((x-A)^2-D^2\bigr),
\]
y su vértice es
\[
V=(A,-aD^2).
\]
Para cada abscisa \(o\), la vertical \(x=o\) corta la gráfica en un único punto
\[
Y_o=\bigl(o,a((o-A)^2-D^2)\bigr).
\]
Es importante decir \emph{el punto de la parábola de abscisa \(o\)}, y no ``el
segundo punto de intersección con la vertical'', porque una gráfica cuadrática tiene
un solo punto sobre cada vertical.

\begin{theorem}[Realización parabólica]\label{thm:parabola-harmonic}
Si \(o\ne A\), la recta \(VY_o\) corta el eje \(x\) en el punto de abscisa
\[
h_{r,s}(o)=A+\frac{D^2}{o-A}.
\]
Cuando \(o=A\), la recta límite es la tangente horizontal en el vértice, cuyo punto
de intersección proyectivo con el eje es \(\infty=h_{r,s}(A)\).
\end{theorem}

\begin{proof}
La pendiente de \(VY_o\) es
\[
\frac{a((o-A)^2-D^2)+aD^2}{o-A}=a(o-A).
\]
Por tanto, la recta tiene ecuación
\[
y+aD^2=a(o-A)(x-A).
\]
Al imponer \(y=0\), obtenemos
\[
x=A+\frac{D^2}{o-A}=h_{r,s}(o).
\]
El caso \(o=A\) se obtiene como límite proyectivo.
\end{proof}

\begin{corollary}[Lectura de dos medias]
Para \(r,s>0\),
\[
\AM(r,s)=h_{r,s}(\infty)=\frac{r+s}{2},
\qquad
\HM(r,s)=h_{r,s}(0)=\frac{2rs}{r+s}.
\]
\end{corollary}

\begin{figure}[p]
\centering
\begin{tikzpicture}
\begin{axis}[
  width=0.92\textwidth,height=0.58\textwidth,
  axis lines=middle,xmin=-2.7,xmax=5.3,ymin=-3.2,ymax=6.7,
  xlabel={$x$},ylabel={$y$},
  xtick={-2,0,1,1.6,2,2.5,4,5},
  xticklabels={$-2$,$0$,$1$,$\frac85$,$2$,$\frac52$,$4$,$5$},
  ytick={-2.25,0,4},yticklabels={$-\frac94$,$0$,$4$},
  grid=both,minor tick num=1,grid style={SoftGray},
  legend style={draw=none,fill=none,font=\small,at={(0.03,0.97)},anchor=north west},
  clip=false]
\addplot[DeepBlue,very thick,domain=-2.2:5.0,samples=240]{(x-1)*(x-4)};
\addlegendentry{$y=(x-1)(x-4)$}
\addplot[Brick,thick,domain=-2.5:2.55]{-0.9*(x-2.5)-2.25};
\addlegendentry{recta $VY_0$}
\addplot[Green,thick,domain=-2.4:4.1]{-0.5*(x-2.5)-2.25};
\addlegendentry{recta $VY_2$}
\addplot[Slate,dashed] coordinates {(2.5,-3.1) (2.5,6.2)};
\addplot[Slate,dashed] coordinates {(0,-3.1) (0,4.2)};
\addplot[Slate,dashed] coordinates {(2,-3.1) (2,4.2)};
\addplot[only marks,mark=*,mark size=2.7pt,Brick] coordinates {(1,0) (4,0)};
\addplot[only marks,mark=*,mark size=2.7pt,Gold] coordinates {(2.5,-2.25)};
\addplot[only marks,mark=*,mark size=2.7pt,Brick] coordinates {(0,4)};
\addplot[only marks,mark=*,mark size=2.7pt,Green] coordinates {(2,-2)};
\addplot[only marks,mark=*,mark size=2.7pt,Brick] coordinates {(1.6,0)};
\addplot[only marks,mark=*,mark size=2.7pt,Green] coordinates {(-2,0)};
\node[anchor=south east,text=Brick] at (axis cs:1,0) {$R$};
\node[anchor=south west,text=Brick] at (axis cs:4,0) {$S$};
\node[anchor=north west,text=Gold] at (axis cs:2.5,-2.25) {$V$};
\node[anchor=south west,text=Brick] at (axis cs:0,4) {$Y_0$};
\node[anchor=north east,text=Green] at (axis cs:2,-2) {$Y_2$};
\node[anchor=south,text=Brick] at (axis cs:1.6,0) {$H_0=\HM=\frac85$};
\node[anchor=south,text=Green] at (axis cs:-2,0) {$H_2=h_{1,4}(2)$};
\node[anchor=south west,text=Gold] at (axis cs:2.5,0) {$\AM=\frac52$};
\end{axis}
\end{tikzpicture}
\caption{La parábola realiza toda la involución armónica, no sólo los valores especiales \(0\) y \(\infty\). En el ejemplo \(r=1\), \(s=4\), la secante \(VY_0\) corta el eje en \(8/5\), mientras que \(VY_2\) lo corta en \(-2\).}
\label{fig:parabola}
\end{figure}

\subsection{La media geométrica pertenece a otra involución}

La transformación recíproca
\[
m_{r,s}(x)=\frac{rs}{x}
\]
intercambia \(r\) y \(s\), y sus puntos fijos satisfacen \(x^2=rs\). Por tanto,
para \(r,s>0\), el punto fijo positivo es \(\GM(r,s)\). Las involuciones
\(h_{r,s}\) y \(m_{r,s}\) no son iguales: la primera fija las raíces; la segunda las
intercambia. Su compatibilidad se expresa mediante
\[
m_{r,s}(\AM)=\HM,
\qquad
\AM\cdot\HM=\GM^2.
\]
Más adelante, la identidad de producto se generalizará a una cadena completa:
\[
\cN_1\cN_2\cdots\cN_n=\GM^n.
\]


% =============================================================================
\section{Polinomios simétricos, log-concavidad y una cadena de medias}
% =============================================================================

\subsection{De las raíces a los coeficientes}

Sea \(n\ge2\) y sean \(r_1,\dots,r_n\) números, inicialmente en un cuerpo de
característica cero. Los polinomios simétricos elementales son
\[
e_0=1,
\qquad
e_k=\sum_{1\le i_1<\cdots<i_k\le n}
 r_{i_1}\cdots r_{i_k}
\quad(1\le k\le n).
\]
El polinomio mónico de raíces \(r_i\) es
\begin{equation}\label{eq:P-viete}
P(t)=\prod_{i=1}^{n}(t-r_i)
=\sum_{k=0}^{n}(-1)^ke_k t^{n-k}.
\end{equation}
Las identidades de Viète no son sólo un mecanismo de cálculo: convierten un conjunto
de puntos en una lista de coordenadas algebraicas invariantes bajo permutaciones.

Normalizamos esos coeficientes mediante
\begin{equation}\label{eq:ak}
a_k=\frac{e_k}{\binom nk},\qquad 0\le k\le n.
\end{equation}
La normalización es esencial. Si todas las variables son iguales a \(x\), entonces
\[
e_k=\binom nk x^k,
\qquad a_k=x^k.
\]
Por tanto, los cocientes \(a_k/a_{k-1}\) devuelven exactamente \(x\), como debe
hacer una media.

\begin{definition}[Cocientes de Newton normalizados]\label{def:newton-quotient-means}
Para \(r_1,\dots,r_n>0\), definimos
\begin{equation}\label{eq:Nk}
\boxed{
\cN_k(r_1,\dots,r_n)
=\frac{a_k}{a_{k-1}}
=\frac{k}{n-k+1}\frac{e_k}{e_{k-1}}},
\qquad 1\le k\le n.
\end{equation}
Los llamaremos \emph{cocientes de Newton normalizados} o, cuando queramos enfatizar
sus propiedades de internalidad, \emph{medias cociente de Newton}.
\end{definition}

\begin{remark}
La terminología se usa aquí para organizar el manuscrito. Las razones consecutivas de
sumas simétricas elementales aparecen naturalmente en la teoría de desigualdades y de
polinomios reales; no se afirma que la expresión algebraica sea inédita.
\end{remark}

\subsection{Prueba autocontenida de las desigualdades de Newton}

La cadena \(\cN_1\ge\cdots\ge\cN_n\) será consecuencia de la log-concavidad de
\((a_k)\). Como esta propiedad sostiene una parte central del artículo, daremos una
prueba completa.

\begin{theorem}[Newton normalizado]\label{thm:newton}
Si \(r_1,\dots,r_n\in\RR\) y el polinomio
\[
E(x)=\prod_{i=1}^{n}(1+r_i x)
=\sum_{k=0}^{n}\binom nk a_kx^k
\]
tiene todas sus raíces reales, entonces
\begin{equation}\label{eq:newton-logconcave}
a_k^2\ge a_{k-1}a_{k+1},
\qquad 1\le k\le n-1.
\end{equation}
Si todos los \(r_i\) son positivos, la desigualdad es estricta salvo en el caso en que
todos los \(r_i\) sean iguales.
\end{theorem}

\begin{proof}
Fijemos \(k\in\{1,\dots,n-1\}\) y escribamos \(m=n-k+1\). Como la derivación
preserva la realidad de las raíces por el teorema de Rolle, el polinomio
\[
G(x)=\frac{1}{n(n-1)\cdots(n-k+2)}E^{(k-1)}(x)
\]
tiene todas sus raíces reales. Usando
\[
\frac{\dd^{k-1}}{\dd x^{k-1}}\binom n{k-1+j}x^{k-1+j}
=\frac{n!}{(n-k+1)!}\binom mj x^j,
\]
obtenemos, salvo un factor positivo irrelevante,
\[
G(x)=\sum_{j=0}^{m}\binom mj a_{k-1+j}x^j.
\]
El polinomio recíproco
\[
G^*(x)=x^mG(1/x)
=\sum_{j=0}^{m}\binom mj a_{k-1+j}x^{m-j}
\]
también tiene todas sus raíces reales. Derivamos \(m-2\) veces. Sólo los tres
primeros términos contribuyen a los grados \(2,1,0\), y se obtiene
\[
\frac{2}{m!}\frac{\dd^{m-2}}{\dd x^{m-2}}G^*(x)
=a_{k-1}x^2+2a_kx+a_{k+1}.
\]
Este polinomio cuadrático tiene raíces reales; por tanto, su discriminante es no
negativo:
\[
(2a_k)^2-4a_{k-1}a_{k+1}\ge0.
\]
Esto es exactamente \cref{eq:newton-logconcave}.

Para variables positivas, la condición de igualdad en las desigualdades de Newton
implica que las raíces negativas de \(E\) coinciden, es decir, que
\(-1/r_1=\cdots=-1/r_n\). Equivalentemente, \(r_1=\cdots=r_n\).
\end{proof}

\begin{corollary}[Cadena decreciente]\label{cor:chain-monotone}
Para \(r_i>0\),
\begin{equation}\label{eq:chain}
\boxed{
\cN_1\ge\cN_2\ge\cdots\ge\cN_n}.
\end{equation}
La cadena es estricta si no todas las variables son iguales.
\end{corollary}

\begin{proof}
La desigualdad \(a_k^2\ge a_{k-1}a_{k+1}\), con todos los términos positivos, es
equivalente a
\[
\frac{a_k}{a_{k-1}}\ge\frac{a_{k+1}}{a_k}.
\]
El lado izquierdo es \(\cN_k\) y el derecho es \(\cN_{k+1}\).
\end{proof}

\subsection{Los extremos y el producto telescópico}

\begin{theorem}[De la media aritmética a la armónica]\label{thm:endpoints-product}
Para \(r_1,\dots,r_n>0\),
\begin{align}
\cN_1&=\frac{e_1}{n}=\AM(r_1,\dots,r_n),\label{eq:N1-AM}\\
\cN_n&=\frac{ne_n}{e_{n-1}}
=\frac{n}{\sum_{i=1}^{n}r_i^{-1}}
=\HM(r_1,\dots,r_n),\label{eq:Nn-HM}\\
\prod_{k=1}^{n}\cN_k&=e_n=\prod_{i=1}^{n}r_i=\GM(r_1,\dots,r_n)^n.\label{eq:product-GM}
\end{align}
Por tanto,
\begin{equation}\label{eq:main-chain-product}
\boxed{
\AM=\cN_1\ge\cN_2\ge\cdots\ge\cN_n=\HM,
\qquad
\GM=\left(\prod_{k=1}^{n}\cN_k\right)^{1/n}.}
\end{equation}
\end{theorem}

\begin{proof}
La primera identidad es inmediata. Para la segunda, observamos que
\[
e_{n-1}=\sum_{j=1}^{n}\prod_{i\ne j}r_i
=e_n\sum_{j=1}^{n}\frac1{r_j}.
\]
Así,
\[
\frac{ne_n}{e_{n-1}}
=\frac{n}{\sum_j1/r_j}.
\]
Finalmente, el producto de los cocientes es telescópico:
\[
\prod_{k=1}^{n}\cN_k
=\prod_{k=1}^{n}\frac{a_k}{a_{k-1}}
=\frac{a_n}{a_0}=e_n.
\]
\end{proof}

\begin{figure}[ht]
\centering
\begin{tikzpicture}[>=Latex,node distance=9mm]
\foreach \k/\lab/\col in {1/{\AM=\cN_1}/DeepBlue,2/{\cN_2}/Gold,3/{\cN_3}/Gold,4/{\cdots}/Slate,5/{\cN_{n-1}}/Gold,6/{\cN_n=\HM}/Brick}{
  \node[circle,draw=\col,fill=white,minimum size=13mm,font=\small\bfseries] (n\k) at (1.8*\k,0) {$\lab$};
}
\foreach \k in {1,...,5}{
  \pgfmathtruncatemacro{\j}{\k+1}
  \draw[-{Latex[length=2.3mm]},thick,Slate] (n\k)--node[above,font=\scriptsize] {$\ge$}(n\j);
}
\draw[decorate,decoration={brace,mirror,amplitude=6pt},Green,thick]
  ($(n1.south west)+(0,-.35)$)--($(n6.south east)+(0,-.35)$)
  node[midway,below=9pt,text=Green,font=\small\bfseries]
  {$\displaystyle \cN_1\cN_2\cdots\cN_n=\GM^n$};
\end{tikzpicture}
\caption{La media geométrica no es, en general, uno de los peldaños: es la media geométrica de todos ellos.}
\label{fig:newton-chain}
\end{figure}

\subsection[Por qué cada cociente merece llamarse media]{Por qué cada \(\cN_k\) merece llamarse media}

Una función de varias variables positivas puede llamarse media en un sentido amplio
si es simétrica, continua, homogénea de grado uno, idempotente e interna. Es deseable,
además, que sea creciente en cada variable.

\begin{theorem}[Propiedades de media]\label{thm:mean-properties}
Para cada \(k\in\{1,\dots,n\}\), la función \(\cN_k:(0,\infty)^n\to(0,\infty)\)
es:
\begin{enumerate}[label=\textup{(\roman*)}]
\item simétrica;
\item continua y homogénea de grado uno;
\item idempotente: \(\cN_k(x,\dots,x)=x\);
\item estrictamente creciente en cada variable;
\item interna:
\[
\min_i r_i\le\cN_k(r_1,\dots,r_n)\le\max_i r_i.
\]
\end{enumerate}
\end{theorem}

\begin{proof}
Las propiedades (i)--(iii) se leen directamente de \cref{eq:Nk}. Para la monotonicidad,
fijemos un índice \(j\) y separemos la variable \(x=r_j\). Denotemos por
\(E_q\) el polinomio simétrico elemental de grado \(q\) en las otras \(n-1\)
variables, con la convención \(E_q=0\) fuera del intervalo \(0\le q\le n-1\). Entonces
\[
e_k=E_k+xE_{k-1},
\qquad
e_{k-1}=E_{k-1}+xE_{k-2}.
\]
Para \(k\ge2\),
\[
\frac{\partial}{\partial x}\frac{e_k}{e_{k-1}}
=\frac{E_{k-1}^2-E_kE_{k-2}}{(E_{k-1}+xE_{k-2})^2}.
\]
Las desigualdades de Newton aplicadas a las \(n-1\) variables restantes implican
\(E_{k-1}^2>E_kE_{k-2}\) en el cono positivo. Por tanto, la derivada es positiva.
Para \(k=1\), \(\cN_1=e_1/n\) y la derivada es \(1/n\).

Finalmente, si \(m=\min_i r_i\) y \(M=\max_i r_i\), la monotonicidad coordenada a
coordenada da
\[
\cN_k(m,\dots,m)\le\cN_k(r_1,\dots,r_n)
\le\cN_k(M,\dots,M).
\]
La idempotencia convierte los extremos en \(m\) y \(M\).
\end{proof}

\subsection{Dualidad por inversión}

La relación clásica entre media aritmética y armónica se amplía a toda la cadena.

\begin{theorem}[Dualidad recíproca]\label{thm:reciprocal-duality}
Para \(r_i>0\),
\begin{equation}\label{eq:reciprocal-duality}
\boxed{
\cN_k(r_1^{-1},\dots,r_n^{-1})
=\frac{1}{\cN_{n-k+1}(r_1,\dots,r_n)}}.
\end{equation}
\end{theorem}

\begin{proof}
La identidad elemental
\[
e_k(r_1^{-1},\dots,r_n^{-1})
=\frac{e_{n-k}(r_1,\dots,r_n)}{e_n(r_1,\dots,r_n)}
\]
implica, usando \(\binom nk=\binom n{n-k}\), que
\[
a_k(r^{-1})=\frac{a_{n-k}(r)}{a_n(r)}.
\]
Por tanto,
\[
\cN_k(r^{-1})
=\frac{a_{n-k}/a_n}{a_{n-k+1}/a_n}
=\frac{a_{n-k}}{a_{n-k+1}}
=\frac{1}{\cN_{n-k+1}(r)}.
\]
\end{proof}

\begin{keybox}[Qué se ha obtenido sin geometría]
Hasta aquí, la cadena completa se deduce de álgebra simétrica y de las desigualdades
de Newton. Todavía falta responder la pregunta estructural: ¿por qué los números
\(a_k\) deben aparecer como coordenadas de un único punto geométrico y por qué los
cocientes consecutivos corresponden a proyecciones naturales? Esa es la función de la
curva racional normal y de la apolaridad.
\end{keybox}

% =============================================================================
\section{Curva racional normal y espacios osculadores}
% =============================================================================

\subsection{La aplicación de Veronese}

Sea \(\KK\) un cuerpo y supongamos \(n\ge2\). Definimos
\begin{equation}\label{eq:veronese}
\nu_n:\PP^1(\KK)\longrightarrow\PP^n(\KK),
\qquad
[u:v]\longmapsto[u^n:u^{n-1}v:\cdots:uv^{n-1}:v^n].
\end{equation}
Su imagen se llama \emph{curva racional normal de grado \(n\)}. En la carta
\(v\ne0\), con \(t=u/v\), escribimos
\[
\nu_n(t)=[t^n:t^{n-1}:\cdots:t:1].
\]
Los puntos distinguidos son
\[
\nu_n(0)=[0:\cdots:0:1],
\qquad
\nu_n(\infty)=[1:0:\cdots:0].
\]

La palabra ``normal'' indica que la curva de grado \(n\) ocupa el espacio proyectivo
mínimo posible \(\PP^n\) sin quedar contenida en un hiperplano. No debe confundirse
con una noción métrica de normalidad.

\subsection{Qué significa oscilar una curva}

En cálculo, la recta tangente aproxima una curva hasta primer orden. En geometría
proyectiva superior se consideran aproximaciones de orden mayor.

\begin{definition}[Espacio osculador]
Para \(r\in\KK\) y \(0\le m\le n\), definimos
\[
\Osc_r^m(\nu_n)
=\PP\left(\Span\{\nu_n(r),\nu_n'(r),\dots,\nu_n^{(m)}(r)\}\right).
\]
Es un subespacio proyectivo de dimensión \(m\), siempre que la característica no
produzca degeneraciones. El caso \(m=1\) es la tangente y el caso \(m=n-1\) es el
hiperplano osculador.
\end{definition}

\begin{lemma}[Banderas opuestas en \(0\) y \(\infty\)]\label{lem:flags}
En coordenadas \([X_0:\cdots:X_n]\),
\begin{align}
\Osc_0^m(\nu_n)&=\PP\Span\{e_{n-m},\dots,e_n\},\label{eq:osc0}\\
\Osc_\infty^m(\nu_n)&=\PP\Span\{e_0,\dots,e_m\},\label{eq:oscinf}
\end{align}
donde \(e_0,\dots,e_n\) es la base canónica de \(\KK^{n+1}\).
\end{lemma}

\begin{proof}
En la carta afín,
\[
\nu_n(t)=(t^n,t^{n-1},\dots,t,1).
\]
Al evaluar las derivadas de órdenes \(0,\dots,m\) en \(t=0\), sólo aparecen los
últimos \(m+1\) vectores de la base. Para \(\infty\), usamos el parámetro local
\(s=v/u\):
\[
\nu_n([1:s])=(1,s,s^2,\dots,s^n),
\]
y las derivadas en \(s=0\) generan los primeros \(m+1\) vectores.
\end{proof}

\begin{figure}[ht]
\centering
\begin{tikzpicture}[>=Latex,scale=0.92]
\draw[rounded corners,fill=Ice,draw=Navy] (0,0) rectangle (14,4.3);
\foreach \i in {0,...,8}{
  \node[circle,draw=DeepBlue,fill=white,minimum size=8mm] (e\i) at (1.1+1.45*\i,2.15) {$e_{\i}$};
}
\draw[decorate,decoration={brace,amplitude=7pt},Brick,thick]
  ($(e0.north west)+(0,.15)$)--($(e3.north east)+(0,.15)$)
  node[midway,above=8pt,font=\small\bfseries,text=Brick] {$\Osc_\infty^3$};
\draw[decorate,decoration={brace,mirror,amplitude=7pt},Green,thick]
  ($(e5.south west)+(0,-.15)$)--($(e8.south east)+(0,-.15)$)
  node[midway,below=8pt,font=\small\bfseries,text=Green] {$\Osc_0^3$};
\draw[-{Latex[length=2.5mm]},Slate,thick] (e0.west)+(-.7,0)--(e0.west)
 node[left,font=\small] {$\infty$};
\draw[-{Latex[length=2.5mm]},Slate,thick] (e8.east)+(.7,0)--(e8.east)
 node[right,font=\small] {$0$};
\node[font=\small\sffamily,text=Slate] at (7,0.55)
 {dos banderas osculadoras que avanzan desde extremos opuestos};
\end{tikzpicture}
\caption{Las banderas osculadoras en \(\infty\) y \(0\) seleccionan bloques consecutivos de coordenadas. Esta estructura permitirá leer todos los cocientes \(a_k/a_{k-1}\).}
\label{fig:opposite-flags}
\end{figure}

% =============================================================================
\section{Apolaridad y el punto que comprime todas las raíces}
% =============================================================================

\subsection{La forma apolar}

A partir de ahora suponemos
\[
\operatorname{char}(\KK)=0
\quad\text{o bien}\quad
\operatorname{char}(\KK)>n.
\]
Esta hipótesis garantiza que los coeficientes binomiales relevantes son invertibles.

\begin{definition}[Emparejamiento apolar]\label{def:apolar-pairing}
Para \(X=(X_0,\dots,X_n)\) y \(Y=(Y_0,\dots,Y_n)\) en \(\KK^{n+1}\), definimos
\begin{equation}\label{eq:apolar}
B_n(X,Y)=\sum_{i=0}^{n}(-1)^i\binom ni X_iY_{n-i}.
\end{equation}
\end{definition}

Su matriz es antidiagonal, con entradas no nulas
\[
1,-n,\binom n2,\dots,(-1)^n.
\]
Por la hipótesis de característica, ninguna se anula; por tanto, \(B_n\) es no
degenerada. Además,
\[
B_n(Y,X)=(-1)^nB_n(X,Y).
\]
Para nuestros propósitos, la ecuación de ortogonalidad \(B_n(X,Y)=0\) es simétrica,
aunque la forma sea alternante cuando \(n\) es impar.

\begin{proposition}[Primera identidad maestra]\label{prop:first-master}
Para cualesquiera \(t,r\in\KK\),
\begin{equation}\label{eq:first-master}
B_n(\nu_n(t),\nu_n(r))=(t-r)^n.
\end{equation}
En consecuencia, el hiperplano polar de \(\nu_n(r)\),
\[
H_r=\{X\in\PP^n:B_n(X,\nu_n(r))=0\},
\]
coincide con el hiperplano osculador \(\Osc_r^{n-1}(\nu_n)\).
\end{proposition}

\begin{proof}
Sustituyendo \(\nu_n(t)=[t^n:t^{n-1}:\cdots:1]\) y
\(\nu_n(r)=[r^n:r^{n-1}:\cdots:1]\), obtenemos
\[
B_n(\nu_n(t),\nu_n(r))
=\sum_{i=0}^{n}(-1)^i\binom ni t^{n-i}r^i
=(t-r)^n.
\]
Al derivar respecto de \(t\) hasta orden \(n-1\) y evaluar en \(t=r\), se obtiene
\[
B_n(\nu_n^{(j)}(r),\nu_n(r))=0,
\qquad 0\le j\le n-1.
\]
Por tanto, el hiperplano polar contiene el espacio osculador de dimensión \(n-1\), y
ambos deben coincidir.
\end{proof}

\subsection{El punto apolar asociado a un polinomio}

Sea \(P\) como en \cref{eq:P-viete}. Definimos
\begin{equation}\label{eq:KP}
K_P=
\left[
 e_n:\frac{e_{n-1}}{n}:\frac{e_{n-2}}{\binom n2}:\cdots:
 \frac{e_1}{n}:1
\right]
=[a_n:a_{n-1}:\cdots:a_1:a_0].
\end{equation}
La inversión del orden no es caprichosa: está adaptada a la antidiagonal de la forma
apolar.

\begin{proposition}[Segunda identidad maestra]\label{prop:second-master}
Para todo \(t\),
\begin{equation}\label{eq:second-master}
B_n(\nu_n(t),K_P)=P(t).
\end{equation}
En consecuencia, el hiperplano polar \(K_P^\perp\) contiene los puntos
\(\nu_n(r_1),\dots,\nu_n(r_n)\). Si las raíces son distintas, esos \(n\) puntos
son linealmente independientes y \(K_P^\perp\) es el único hiperplano que los
contiene: el hiperplano \(n\)-secante determinado por el divisor de raíces.
\end{proposition}

\begin{proof}
La coordenada de índice \(n-i\) de \(K_P\) es \(e_i/\binom ni\). Por tanto,
\[
B_n(\nu_n(t),K_P)
=\sum_{i=0}^{n}(-1)^i\binom ni t^{n-i}\frac{e_i}{\binom ni}
=\sum_{i=0}^{n}(-1)^ie_it^{n-i}=P(t).
\]
Si \(t=r_j\), el lado derecho se anula, de modo que \(\nu_n(r_j)\) pertenece al
hiperplano polar de \(K_P\). Para raíces distintas, el menor formado por las últimas \(n\) coordenadas de
los vectores \(\nu_n(r_j)\) es la matriz de Vandermonde
\(\bigl(r_j^{\,n-1-i}\bigr)_{0\le i\le n-1,\,1\le j\le n}\), cuyo
determinante es no nulo. Por tanto, los puntos
abarcan un hiperplano proyectivo único, que necesariamente es \(K_P^\perp\).
\end{proof}

\begin{theorem}[Compresión apolar de las raíces]\label{thm:KP-intersection}
Supongamos que \(r_1,\dots,r_n\) son distintas. Entonces
\begin{equation}\label{eq:KP-intersection}
\boxed{
K_P=\bigcap_{i=1}^{n}\Osc_{r_i}^{n-1}(\nu_n).}
\end{equation}
Es decir, el mismo punto que es polar al hiperplano secante determinado por las raíces
es la intersección de los hiperplanos osculadores en esas raíces.
\end{theorem}

\begin{proof}
Sea \(X\) un punto perteneciente a todos los hiperplanos osculadores. Por
\cref{prop:first-master},
\[
B_n(X,\nu_n(r_i))=0
\]
para cada \(i\). Como la ortogonalidad es simétrica hasta signo,
\(B_n(\nu_n(r_i),X)=0\). Definamos
\[
F_X(t)=B_n(\nu_n(t),X).
\]
Es un polinomio de grado a lo sumo \(n\). La no degeneración de \(B_n\) garantiza
que no es idénticamente nulo para un punto proyectivo \(X\). Como tiene las \(n\)
raíces distintas \(r_i\), debe ser proporcional a \(P\):
\[
F_X(t)=\lambda P(t).
\]
Pero \cref{prop:second-master} afirma que \(P(t)=B_n(\nu_n(t),K_P)\). Comparando
coeficientes se obtiene \(X=\lambda K_P\), es decir, el mismo punto proyectivo.
\end{proof}

\begin{figure}[ht]
\centering
\begin{tikzpicture}[>=Latex,node distance=10mm and 14mm,scale=0.76,transform shape]
\node[draw=Navy,fill=Ice,rounded corners,minimum width=42mm,minimum height=13mm] (roots)
 {$r_1,\dots,r_n\in\PP^1$};
\node[draw=DeepBlue,fill=Sky,rounded corners,minimum width=52mm,minimum height=15mm,right=of roots] (points)
 {$\nu_n(r_1),\dots,\nu_n(r_n)\in\PP^n$};
\node[draw=Brick,fill=white,rounded corners,minimum width=46mm,minimum height=15mm,above right=7mm and 15mm of points] (sec)
 {hiperplano secante $\Pi_P$};
\node[draw=Green,fill=Mint,rounded corners,minimum width=49mm,minimum height=15mm,below right=7mm and 15mm of points] (osc)
 {$n$ hiperplanos osculadores};
\node[draw=Gold,fill=Warm,rounded corners,minimum width=48mm,minimum height=17mm,right=35mm of points] (kp)
 {$K_P=[a_n:\cdots:a_0]$};
\draw[->,thick,DeepBlue] (roots)--node[above,font=\small] {$\nu_n$}(points);
\draw[->,thick,Brick] (points)--(sec);
\draw[->,thick,Green] (points)--(osc);
\draw[<->,thick,Gold] (sec)--node[above,sloped,font=\small] {polaridad}(kp);
\draw[->,thick,Green] (osc)--node[below,sloped,font=\small] {intersección}(kp);
\node[font=\small\sffamily,text=Slate,align=center] at ($(points.south)+(0,-2.1)$)
 {dos descripciones duales del mismo punto};
\end{tikzpicture}
\caption{El punto \(K_P\) comprime toda la información simétrica de las raíces.}
\label{fig:apolar-compression}
\end{figure}


% =============================================================================
\section{Banderas osculadoras opuestas y realización geométrica de toda la cadena}
% =============================================================================

\subsection{Cómo seleccionar dos coordenadas consecutivas}

El punto \(K_P=[a_n:a_{n-1}:\cdots:a_0]\) contiene toda la cadena, pero aún debemos
explicar por qué cada razón \(a_k/a_{k-1}\) es una proyección geométrica natural.
Las banderas de \cref{lem:flags} proporcionan exactamente los centros y las rectas
objetivo necesarios.

Para \(1\le k\le n\), definimos la recta coordenada
\begin{equation}\label{eq:Lk}
L_k=\PP\Span\{e_{n-k},e_{n-k+1}\}
\end{equation}
y el subespacio complementario
\begin{equation}\label{eq:Lambdak}
\Lambda_k=
\PP\Span\{e_0,\dots,e_{n-k-1},e_{n-k+2},\dots,e_n\}.
\end{equation}
Usaremos la convención de que un bloque vacío simplemente se omite. Por ejemplo,
\[
\Lambda_1=\PP\Span\{e_0,\dots,e_{n-2}\}=\Osc_\infty^{n-2}(\nu_n),
\]
\[
\Lambda_n=\PP\Span\{e_2,\dots,e_n\}=\Osc_0^{n-2}(\nu_n).
\]
Para los índices intermedios,
\begin{equation}\label{eq:center-flags}
\Lambda_k=
\left\langle
\Osc_\infty^{n-k-1}(\nu_n),
\Osc_0^{k-2}(\nu_n)
\right\rangle.
\end{equation}
Es la unión lineal de dos piezas tomadas de banderas osculadoras opuestas.

Como los espacios vectoriales que subyacen a \(L_k\) y \(\Lambda_k\) son
complementarios, existe una proyección lineal
\[
\pi_k:\PP^n\dashrightarrow L_k
\]
con centro \(\Lambda_k\). En coordenadas, esta proyección simplemente conserva
\(X_{n-k}\) y \(X_{n-k+1}\):
\begin{equation}\label{eq:pik}
\pi_k([X_0:\cdots:X_n])
=[0:\cdots:0:X_{n-k}:X_{n-k+1}:0:\cdots:0].
\end{equation}

\begin{theorem}[Realización proyectiva de la cadena]\label{thm:projective-chain}
Para \(K_P=[a_n:a_{n-1}:\cdots:a_0]\),
\begin{equation}\label{eq:pikKP}
\pi_k(K_P)=[a_k:a_{k-1}]\in L_k.
\end{equation}
Si identificamos \(L_k\simeq\PP^1\) mediante la coordenada afín
\(z=X_{n-k}/X_{n-k+1}\), entonces
\begin{equation}\label{eq:pik-coordinate}
z(\pi_k(K_P))=\frac{a_k}{a_{k-1}}=\cN_k.
\end{equation}
En particular, los dos extremos son las proyecciones desde los osculadores de orden
\(n-2\):
\begin{align}
\pi_1(K_P)&=[a_1:a_0] &&\Longrightarrow&& z=\AM,\label{eq:proj-AM}\\
\pi_n(K_P)&=[a_n:a_{n-1}] &&\Longrightarrow&& z=\HM.\label{eq:proj-HM}
\end{align}
\end{theorem}

\begin{proof}
La forma de \(K_P\) en \cref{eq:KP} muestra que su coordenada de índice \(n-k\) es
\(a_k\) y la siguiente es \(a_{k-1}\). La proyección \(\pi_k\) elimina todas las
demás coordenadas. La razón afín es, por definición, \(a_k/a_{k-1}\).
\end{proof}

\begin{figure}[p]
\centering
\begin{tikzpicture}[>=Latex,scale=0.94]
\draw[rounded corners,fill=Ice,draw=Navy] (0,0) rectangle (15.2,7.9);
\node[font=\sffamily\bfseries,text=Navy] at (7.6,7.45)
 {El punto apolar y sus ventanas de dos coordenadas};

\foreach \i/\lab in {0/{a_n},1/{a_{n-1}},2/{a_{n-2}},3/{\cdots},4/{a_k},5/{a_{k-1}},6/{\cdots},7/{a_1},8/{a_0}}{
  \node[draw=DeepBlue,fill=white,rounded corners,minimum width=12mm,minimum height=10mm] (a\i) at (1.1+1.6*\i,5.65) {$\lab$};
}
\node[left=2mm of a0,font=\small\bfseries] {$K_P=$};
\draw[decorate,decoration={brace,amplitude=7pt},Green,thick]
  ($(a0.north west)+(0,.15)$)--($(a3.north east)+(0,.15)$)
  node[midway,above=8pt,font=\small,text=Green] {$\Osc_\infty^{n-k-1}$};
\draw[decorate,decoration={brace,amplitude=7pt},Brick,thick]
  ($(a6.south west)+(0,-.15)$)--($(a8.south east)+(0,-.15)$)
  node[midway,below=8pt,font=\small,text=Brick] {$\Osc_0^{k-2}$};
\draw[very thick,Gold,rounded corners] ($(a4.north west)+(-.13,.13)$) rectangle ($(a5.south east)+(.13,-.13)$);
\node[draw=Gold,fill=Warm,rounded corners,minimum width=34mm,minimum height=13mm] (ratio) at (7.6,2.75)
 {$\displaystyle [a_k:a_{k-1}]$};
\draw[-{Latex[length=3mm]},Gold,line width=1.4pt] ($(a4.south)!0.5!(a5.south)$)--(ratio.north);
\node[draw=Navy,fill=white,rounded corners,minimum width=36mm,minimum height=13mm] (mean) at (7.6,1.0)
 {$\displaystyle \cN_k=\frac{a_k}{a_{k-1}}$};
\draw[-{Latex[length=2.5mm]},Navy,thick] (ratio)--(mean);
\node[font=\small\sffamily,text=Slate,align=center] at (2.2,1.15)
 {el centro $\Lambda_k$\\elimina los bloques exteriores};
\node[font=\small\sffamily,text=Slate,align=center] at (13.0,1.15)
 {la recta $L_k$\\retiene una ventana consecutiva};
\draw[Slate,dashed] (2.2,1.75)--(3.8,4.75);
\draw[Slate,dashed] (13.0,1.75)--(10.1,4.75);
\end{tikzpicture}
\caption{Cada media \(\cN_k\) es la coordenada afín de una proyección de \(K_P\) determinada por las banderas osculadoras opuestas en \(0\) y \(\infty\).}
\label{fig:coordinate-window}
\end{figure}

\subsection{Qué es intrínseco y qué depende de un marco afín}

La construcción de \(K_P\) es proyectivamente natural: cambia covariantemente cuando
se aplica una transformación de \(\PGL_2\) a las raíces y la representación de grado
\(n\) al espacio ambiente. Sin embargo, las cantidades numéricas
\(\AM,\HM,\cN_k\) requieren distinguir \(0\) y \(\infty\). En otras palabras, dependen
de un marco afín sobre \(\PP^1\).

\begin{warning}[No confundir naturalidad con invariancia absoluta]
La media armónica no es invariante bajo una transformación proyectiva arbitraria.
La afirmación correcta es que la construcción es natural para una recta proyectiva
\emph{enmarcada}, es decir, para el dato
\[
(\PP^1;0,\infty;r_1,\dots,r_n).
\]
Al cambiar el marco, cambian las coordenadas y las medias. Lo que permanece es la
arquitectura apolar y osculadora que produce las nuevas coordenadas en el nuevo marco.
\end{warning}

\subsection{La cadena como espectro de cocientes del punto apolar}

El teorema anterior permite pensar en
\[
\cN(P)=\bigl(\cN_1(P),\dots,\cN_n(P)\bigr)
\]
como un espectro de lecturas consecutivas de \(K_P\). La desigualdad de Newton
significa que, para raíces positivas, las coordenadas de \(K_P\) forman una sucesión
log-cóncava:
\[
a_k^2\ge a_{k-1}a_{k+1}.
\]
Geométricamente, la pendiente afín de cada ventana no aumenta al avanzar desde el
extremo de la media aritmética hacia el extremo de la media armónica.

Esta lectura conecta tres lenguajes:
\[
\begin{array}{ccl}
\text{coeficientes de }P &:& a_k=e_k/\binom nk,\\[1mm]
\text{geometría de }K_P &:& \pi_k(K_P)=[a_k:a_{k-1}],\\[1mm]
\text{desigualdades} &:& \cN_k=a_k/a_{k-1}\ge a_{k+1}/a_k.
\end{array}
\]

% =============================================================================
\section{La torre de derivadas normalizadas}
% =============================================================================

\subsection{Una secuencia de polinomios mónicos de todos los grados}

La geometría apolar explica dónde viven las cantidades \(\cN_k\). Ahora construiremos
una segunda ruta, puramente polinomial al principio, que desembocará en dinámica
racional.

\begin{definition}[Torre derivada normalizada]\label{def:Qk}
Sea \(P\) mónico de grado \(n\). Para \(0\le k\le n\), definimos
\begin{equation}\label{eq:Qk}
Q_k(t)=\frac{k!}{n!}P^{(n-k)}(t).
\end{equation}
En particular, \(Q_n=P\) y \(Q_0=1\).
\end{definition}

\begin{proposition}[Estructura de Appell truncada]\label{prop:Qk-properties}
Para \(1\le k\le n\):
\begin{enumerate}[label=\textup{(\roman*)}]
\item \(Q_k\) es mónico de grado \(k\);
\item \(Q_k'=kQ_{k-1}\);
\item su término constante es
\begin{equation}\label{eq:Qk-zero}
Q_k(0)=(-1)^k\frac{e_k}{\binom nk}=(-1)^ka_k;
\end{equation}
\item el coeficiente de \(t^{k-1}\) es \(-k e_1/n\), de modo que la media
aritmética de sus raíces es \(e_1/n\), independiente de \(k\).
\end{enumerate}
\end{proposition}

\begin{proof}
Como \(P\) es mónico, su derivada de orden \(n-k\) tiene término principal
\(n!/k!\,t^k\), y la normalización lo vuelve mónico. Además,
\[
Q_k' = \frac{k!}{n!}P^{(n-k+1)}
= k\frac{(k-1)!}{n!}P^{(n-k+1)}=kQ_{k-1}.
\]
El término constante se obtiene de \cref{eq:P-viete}:
\[
P^{(n-k)}(0)=(n-k)!(-1)^ke_k,
\]
y por tanto
\[
Q_k(0)=\frac{k!(n-k)!}{n!}(-1)^ke_k
=(-1)^k\frac{e_k}{\binom nk}.
\]
Para el coeficiente de \(t^{k-1}\), derivamos el término
\(-e_1t^{n-1}\):
\[
\frac{k!}{n!}\frac{(n-1)!}{(k-1)!}(-e_1)t^{k-1}
=-\frac{k}{n}e_1t^{k-1}.
\]
La suma de las raíces de un polinomio mónico de grado \(k\) es, por Viète,
\(ke_1/n\), y su promedio es \(e_1/n\).
\end{proof}

\begin{figure}[ht]
\centering
\begin{tikzpicture}[>=Latex,node distance=7mm and 8mm,scale=0.72,transform shape]
\node[draw=Navy,fill=Ice,rounded corners,minimum width=29mm,minimum height=12mm] (qn) {$Q_n=P$};
\node[draw=DeepBlue,fill=Sky,rounded corners,minimum width=29mm,minimum height=12mm,right=of qn] (qn1) {$Q_{n-1}=P'/n$};
\node[draw=DeepBlue,fill=Sky,rounded corners,minimum width=29mm,minimum height=12mm,right=of qn1] (qn2) {$Q_{n-2}$};
\node[right=of qn2,font=\Large] (dots) {$\cdots$};
\node[draw=Gold,fill=Warm,rounded corners,minimum width=29mm,minimum height=12mm,right=of dots] (q2) {$Q_2$};
\node[draw=Brick,fill=white,rounded corners,minimum width=29mm,minimum height=12mm,right=of q2] (q1) {$Q_1=t-\AM$};
\foreach \a/\b/\lab in {qn/qn1/{\frac1n\,D},qn1/qn2/{\frac1{n-1}\,D},qn2/dots/{},dots/q2/{},q2/q1/{\frac12\,D}}{
  \draw[-{Latex[length=2.3mm]},thick,Green] (\a)--node[above,font=\scriptsize] {$\lab$}(\b);
}
\node[below=9mm of qn,font=\small\sffamily,text=Slate] {grado $n$};
\node[below=9mm of q1,font=\small\sffamily,text=Slate] {grado $1$};
\draw[decorate,decoration={brace,mirror,amplitude=6pt},Slate]
  ($(qn.south west)+(0,-1.0)$)--($(q1.south east)+(0,-1.0)$)
  node[midway,below=8pt,font=\small\bfseries,text=Slate]
  {todas las capas conservan el mismo centroide de raíces};
\end{tikzpicture}
\caption{La normalización convierte las derivadas sucesivas en una torre de polinomios mónicos con la relación exacta \(Q_k'=kQ_{k-1}\).}
\label{fig:derivative-tower}
\end{figure}

\subsection{Raíces derivadas e intercalación}

Si \(P\) tiene raíces reales, Rolle garantiza que \(P'\) también las tiene y que se
intercalan. Repitiendo el argumento, todos los \(Q_k\) son reales enraizados. Si las
raíces originales son positivas, cada raíz derivada queda en el intervalo convexo
\([\min r_i,\max r_i]\), por lo que también es positiva.

Denotemos las raíces de \(Q_k\) por
\[
\rho_{k,1},\dots,\rho_{k,k}.
\]
Para raíces simples ordenadas, la intercalación toma la forma
\[
r_1<\rho_{n-1,1}<r_2<\rho_{n-1,2}<\cdots<\rho_{n-1,n-1}<r_n.
\]
En niveles inferiores se obtiene una malla triangular de raíces.

\begin{figure}[ht]
\centering
\begin{tikzpicture}[x=1.1cm,y=0.78cm,>=Latex]
\foreach \row/\ys/\pts/\col/\lab in {
0/0/{0.4,2.0,3.6,5.2,7.1}/Navy/{$Q_5=P$},
1/-1/{1.2,2.8,4.45,6.15}/DeepBlue/{$Q_4$},
2/-2/{2.0,3.65,5.35}/Green/{$Q_3$},
3/-3/{2.85,4.5}/Gold/{$Q_2$},
4/-4/{3.7}/Brick/{$Q_1$}}
{
  \draw[Slate!55] (0,\ys)--(7.6,\ys);
  \foreach \x in \pts {\fill[\col] (\x,\ys) circle (2.4pt);}
  \node[right,font=\small\bfseries,text=\col] at (7.8,\ys) {\lab};
}
\foreach \x in {0.4,2.0,3.6,5.2,7.1}{\draw[Slate!25,dashed] (\x,0)--(3.7,-4);}
\node[align=center,font=\small\sffamily,text=Slate] at (3.8,-5.0)
 {Rolle comprime progresivamente la configuración\\sin mover su media aritmética};
\end{tikzpicture}
\caption{Esquema de intercalación para un polinomio real de grado cinco. Las posiciones son ilustrativas; la propiedad esencial es que cada fila se intercala con la anterior.}
\label{fig:interlacing}
\end{figure}

% =============================================================================
\section{La torre polar de medias}
% =============================================================================

\subsection{Un mapa racional en cada nivel}

La relación \(Q_k'=kQ_{k-1}\) sugiere comparar dos niveles consecutivos mediante una
razón racional.

\begin{definition}[Mapa polar del nivel \(k\)]\label{def:Phik}
Para \(1\le k\le n\), definimos
\begin{equation}\label{eq:Phik}
\Phi_k(t)
=t-k\frac{Q_k(t)}{Q_k'(t)}
=t-\frac{Q_k(t)}{Q_{k-1}(t)}.
\end{equation}
Para \(k=1\), interpretamos \(Q_0=1\), de modo que
\[
\Phi_1(t)=t-Q_1(t)=\AM(r_1,\dots,r_n),
\]
un mapa constante.
\end{definition}

La forma \(t-kQ_k/Q_k'\) muestra que \(\Phi_k\) es el mapa polar asociado al
polinomio de grado \(k\), \(Q_k\). La forma \(t-Q_k/Q_{k-1}\) hace visible la
naturaleza telescópica de la torre.

\begin{theorem}[Teorema de la torre polar de medias]\label{thm:polar-tower-means}
Sea \(P(t)=\prod_{i=1}^{n}(t-r_i)\) con \(r_i>0\), y sean \(Q_k\), \(\Phi_k\) y
\(\cN_k\) como antes. Entonces, para cada \(1\le k\le n\):
\begin{enumerate}[label=\textup{(\roman*)}]
\item
\begin{equation}\label{eq:Nk-Phik0}
\boxed{\cN_k(r_1,\dots,r_n)=\Phi_k(0)};
\end{equation}
\item si \(\rho_{k,1},\dots,\rho_{k,k}\) son las raíces de \(Q_k\),
\begin{equation}\label{eq:Nk-HM-derived}
\boxed{
\cN_k(r_1,\dots,r_n)
=\HM(\rho_{k,1},\dots,\rho_{k,k})};
\end{equation}
\item todas las capas tienen la misma media aritmética:
\begin{equation}\label{eq:AM-preserved}
\AM(\rho_{k,1},\dots,\rho_{k,k})
=\AM(r_1,\dots,r_n);
\end{equation}
\item la media geométrica original satisface
\begin{equation}\label{eq:GM-product-derived-HM}
\boxed{
\GM(r_1,\dots,r_n)^n
=\prod_{k=1}^{n}\HM(\rho_{k,1},\dots,\rho_{k,k})}.
\end{equation}
\end{enumerate}
\end{theorem}

\begin{proof}
Por \cref{eq:Qk-zero},
\[
\Phi_k(0)=-\frac{Q_k(0)}{Q_{k-1}(0)}
=-\frac{(-1)^ka_k}{(-1)^{k-1}a_{k-1}}
=\frac{a_k}{a_{k-1}}=\cN_k.
\]
Como \(Q_k\) es mónico de grado \(k\) y tiene raíces positivas \(\rho_{k,j}\), la
identidad general
\[
-k\frac{Q_k(0)}{Q_k'(0)}
=\frac{k}{\sum_{j=1}^{k}\rho_{k,j}^{-1}}
\]
muestra que \(\Phi_k(0)\) es la media armónica de esas raíces. La conservación de la
media aritmética fue probada en \cref{prop:Qk-properties}. Finalmente,
\cref{eq:GM-product-derived-HM} se obtiene sustituyendo
\(\cN_k=\HM(\rho_{k,1},\dots,\rho_{k,k})\) en
\cref{eq:product-GM}.
\end{proof}

\begin{keybox}[Interpretación]
La cadena \(\AM=\cN_1\ge\cdots\ge\cN_n=\HM\) puede leerse como sigue:
\begin{enumerate}[leftmargin=8mm]
\item se deriva \(P\) hasta obtener un polinomio mónico de grado \(k\);
\item se toman sus \(k\) raíces, que siguen teniendo la misma media aritmética;
\item se calcula la media armónica de esas raíces;
\item al aumentar \(k\), esa media armónica desciende desde el centroide único de
\(Q_1\) hasta la media armónica de todas las raíces originales.
\end{enumerate}
\end{keybox}

\subsection{Una fórmula local alrededor de cualquier punto}

La evaluación en \(0\) no es esencial. Si recentramos el polinomio en un punto
\(o\), las raíces relativas son \(r_i-o\). Aplicando la definición a esa lista se
obtiene una familia local.

\begin{definition}[Cociente local de orden \(k\)]
Siempre que el denominador no se anule, definimos
\begin{equation}\label{eq:local-Nk}
\cN_k(P;o)
=-k\frac{P^{(n-k)}(o)}{P^{(n-k+1)}(o)}.
\end{equation}
La cantidad afín trasladada es
\begin{equation}\label{eq:local-map}
\mathfrak M_k(P;o)=o+\cN_k(P;o)
=o-k\frac{P^{(n-k)}(o)}{P^{(n-k+1)}(o)}
=\Phi_k(o).
\end{equation}
\end{definition}

La igualdad con \(\Phi_k\) es inmediata por la normalización de \(Q_k\). En términos
de las raíces relativas,
\[
\cN_k(P;o)=\cN_k(r_1-o,\dots,r_n-o).
\]
Para \(k=1\),
\[
\Phi_1(o)=\AM(r_1,\dots,r_n),
\]
independiente de \(o\). Para \(k=n\),
\[
\Phi_n(o)=o-n\frac{P(o)}{P'(o)},
\]
que será el mapa polar completo.

\subsection[Ejemplo completo: raíces 1, 2 y 3]{Ejemplo completo: raíces \(1,2,3\)}

Tomemos
\[
P(t)=(t-1)(t-2)(t-3)=t^3-6t^2+11t-6.
\]
Entonces
\[
a_0=1,
\qquad a_1=\frac{6}{3}=2,
\qquad a_2=\frac{11}{3},
\qquad a_3=6.
\]
La cadena es
\[
\cN_1=2,
\qquad
\cN_2=\frac{11/3}{2}=\frac{11}{6},
\qquad
\cN_3=\frac{6}{11/3}=\frac{18}{11}.
\]
En efecto,
\[
2>\frac{11}{6}>\frac{18}{11},
\qquad
2\cdot\frac{11}{6}\cdot\frac{18}{11}=6=\GM(1,2,3)^3.
\]

La torre normalizada es
\[
Q_3=P,
\qquad
Q_2=\frac{P'}{3}=t^2-4t+\frac{11}{3},
\qquad
Q_1=\frac{P''}{6}=t-2.
\]
Las raíces de \(Q_2\) son
\[
2\pm\frac{1}{\sqrt3},
\]
y su media armónica es
\[
\frac{2}{\dfrac1{2-1/\sqrt3}+\dfrac1{2+1/\sqrt3}}
=\frac{11}{6}=\cN_2.
\]

\begin{figure}[ht]
\centering
\begin{tikzpicture}
\begin{axis}[
 width=0.92\textwidth,height=0.48\textwidth,
 axis lines=middle,xmin=0.55,xmax=3.45,ymin=-2.7,ymax=3.2,
 xlabel={$t$},ylabel={},
 xtick={1,1.42265,2,2.57735,3},
 xticklabels={$1$,$2-1/\sqrt3$,$2$,$2+1/\sqrt3$,$3$},
 ytick=\empty,grid=major,grid style={SoftGray},
 legend style={draw=none,fill=none,at={(0.02,0.97)},anchor=north west,font=\small}]
\addplot[Navy,very thick,domain=.65:3.35,samples=220]{(x-1)*(x-2)*(x-3)};
\addlegendentry{$Q_3=P$}
\addplot[Green,thick,domain=.65:3.35,samples=220]{x^2-4*x+11/3};
\addlegendentry{$Q_2=P'/3$}
\addplot[Brick,thick,domain=.65:3.35,samples=2]{x-2};
\addlegendentry{$Q_1=P''/6$}
\addplot[only marks,mark=*,mark size=2.4pt,Navy] coordinates {(1,0)(2,0)(3,0)};
\addplot[only marks,mark=*,mark size=2.4pt,Green] coordinates {(1.42265,0)(2.57735,0)};
\addplot[only marks,mark=*,mark size=2.4pt,Brick] coordinates {(2,0)};
\end{axis}
\end{tikzpicture}
\caption{La torre derivada para \(P(t)=(t-1)(t-2)(t-3)\). Las raíces de cada nivel se intercalan y conservan centroide \(2\); sus medias armónicas son \(2,11/6,18/11\).}
\label{fig:cubic-example-tower}
\end{figure}


% =============================================================================
\section{Formas binarias y el mapa polar exterior}
% =============================================================================

\subsection{Homogeneizar el polinomio}

La expresión \(t-kQ_k/Q_k'\) parece depender de una coordenada afín. Para descubrir
su contenido proyectivo, debemos regresar a coordenadas homogéneas.

Sea \(V\) un espacio vectorial de dimensión dos sobre \(\KK\), con coordenadas
\((X,Y)\). Una forma binaria de grado \(m\) es un elemento de
\(\Sym^m(V^*)\), es decir, un polinomio homogéneo
\[
F(X,Y)=\sum_{j=0}^{m}c_jX^{m-j}Y^j.
\]
Si \(Q(t)\) es mónico de grado \(m\), su homogeneización es
\begin{equation}\label{eq:homogenization}
F_Q(X,Y)=Y^mQ(X/Y).
\end{equation}
Para \(Q=P\),
\begin{equation}\label{eq:F-P}
F_P(X,Y)=\prod_{i=1}^{n}(X-r_iY)
=\sum_{k=0}^{n}(-1)^ke_kX^{n-k}Y^k.
\end{equation}
Los ceros de \(F_P\) son puntos de \(\PP^1\), de modo que la forma binaria codifica
el divisor de raíces sin privilegiar una carta afín.

\subsection{Del diferencial a un vector}

Fijemos una forma de área no nula \(\omega\in\bigwedge^2V^*\). En coordenadas,
podemos tomar \(\omega=\dd X\wedge\dd Y\). Para cada punto \(v\in V\), el
diferencial \(\dd F_v\) es una forma lineal. La forma \(\omega\) identifica formas
lineales con vectores, salvo un factor escalar. En coordenadas, el vector hamiltoniano
es proporcional a
\[
\cH_F(X,Y)=(-F_Y(X,Y),F_X(X,Y)).
\]
Como ambas componentes son homogéneas de grado \(m-1\), su proyectivización define
un mapa racional.

\begin{definition}[Mapa polar exterior]\label{def:polar-map}
Sea \(F\) una forma binaria no nula de grado \(m\ge2\). Definimos
\begin{equation}\label{eq:polar-map-homogeneous}
\cD_F:\PP^1\dashrightarrow\PP^1,
\qquad
[X:Y]\longmapsto[-F_Y(X,Y):F_X(X,Y)].
\end{equation}
Cuando \(F_X\) y \(F_Y\) no tienen factor común, el mapa tiene grado \(m-1\).
\end{definition}

El signo global es irrelevante en coordenadas proyectivas. La elección anterior tiene
la ventaja de producir directamente la fórmula afín que necesitamos.

\begin{proposition}[Expresión afín]\label{prop:affine-polar}
Sea \(F_Q(X,Y)=Y^mQ(X/Y)\). En la carta \(Y=1\), el mapa \(\cD_{F_Q}\) es
\begin{equation}\label{eq:PhiQ}
\boxed{
\Phi_Q(t)=t-m\frac{Q(t)}{Q'(t)}}.
\end{equation}
\end{proposition}

\begin{proof}
En \((t,1)\),
\[
(F_Q)_X(t,1)=Q'(t).
\]
La identidad de Euler para una forma homogénea de grado \(m\),
\[
XF_X+YF_Y=mF,
\]
da
\[
(F_Q)_Y(t,1)=mQ(t)-tQ'(t).
\]
Por tanto,
\[
[-(F_Q)_Y:(F_Q)_X]
=[tQ'(t)-mQ(t):Q'(t)],
\]
cuya coordenada afín es \(t-mQ/Q'\).
\end{proof}

Aplicando esta proposición a \(Q_k\), obtenemos exactamente los mapas \(\Phi_k\) de
\cref{def:Phik}. La torre de medias es, por tanto, una torre de mapas polares
exteriores asociados a formas binarias de grados \(1,2,\dots,n\).

\subsection{Covariancia proyectiva}

La construcción no depende esencialmente de las coordenadas.

\begin{theorem}[Covariancia bajo \(\PGL_2\)]\label{thm:covariance}
Sea \(g\in\GL(V)\) y definamos la acción sobre formas por
\[
(g\cdot F)(v)=F(g^{-1}v).
\]
Entonces, como mapas proyectivos,
\begin{equation}\label{eq:covariance}
\boxed{
\cD_{g\cdot F}=g\circ\cD_F\circ g^{-1}.}
\end{equation}
\end{theorem}

\begin{proof}
La regla de la cadena da
\[
\dd(g\cdot F)_v=\dd F_{g^{-1}v}\circ g^{-1}.
\]
La identificación entre \(V^*\) y \(V\) mediante \(\omega\) cambia bajo \(g\) por el
factor \(\det(g)\). Por tanto, el vector hamiltoniano de \(g\cdot F\) en \(v\) es
proporcional a
\[
g\bigl(\cH_F(g^{-1}v)\bigr).
\]
El factor escalar desaparece al proyectivizar, y se obtiene \cref{eq:covariance}.
\end{proof}

\begin{figure}[ht]
\centering
\begin{tikzcd}[column sep=large,row sep=large]
\PP^1 \arrow[r,"{\cD_F}"] \arrow[d,"g"']
& \PP^1 \arrow[d,"g"]\\
\PP^1 \arrow[r,"{\cD_{g\cdot F}}"']
& \PP^1
\end{tikzcd}
\caption{La covariancia significa que transformar la configuración de raíces y después construir el mapa produce el mismo resultado que conjugar el mapa original.}
\label{fig:covariance}
\end{figure}

\subsection{Los extremos de la cadena dentro del mapa polar superior}

Para \(P\) mónico de grado \(n\), el mapa superior es
\begin{equation}\label{eq:PhiP}
\Phi_P(t)=t-n\frac{P(t)}{P'(t)}.
\end{equation}
La forma homogénea permite evaluar sin ambigüedad en \(0\) y \(\infty\).

\begin{proposition}[Dos valores distinguidos]\label{prop:Phi-endpoints}
Si \(F_P\) es como en \cref{eq:F-P}, entonces
\begin{align}
\cD_{F_P}([0:1])&=[ne_n:e_{n-1}],\label{eq:Phi-zero-hom}\\
\cD_{F_P}([1:0])&=[e_1:n].\label{eq:Phi-infinity-hom}
\end{align}
En la carta afín,
\begin{equation}\label{eq:Phi-endpoint-means}
\boxed{
\Phi_P(0)=\HM(r_1,\dots,r_n),
\qquad
\Phi_P(\infty)=\AM(r_1,\dots,r_n).}
\end{equation}
\end{proposition}

\begin{proof}
En \([0:1]\), los únicos términos relevantes son
\[
(F_P)_X(0,1)=(-1)^{n-1}e_{n-1},
\qquad
(F_P)_Y(0,1)=n(-1)^ne_n.
\]
Por tanto,
\[
[-F_Y:F_X]=[ne_n:e_{n-1}].
\]
En \([1:0]\),
\[
(F_P)_X(1,0)=n,
\qquad
(F_P)_Y(1,0)=-e_1,
\]
y el valor es \([e_1:n]\).
\end{proof}

\begin{remark}
Las proyecciones apolares de \cref{thm:projective-chain} y el mapa polar exterior
coinciden en los dos valores extremos, pero no deben confundirse como una única
familia de proyecciones móviles. La primera estructura lee coordenadas de \(K_P\);
la segunda procede del diferencial de la forma binaria.
\end{remark}

% =============================================================================
\section{Puntos fijos, multiplicadores y caracterización por unicidad}
% =============================================================================

\subsection{El divisor de raíces es el divisor de puntos fijos}

La identidad de Euler produce una caracterización extraordinariamente compacta.

\begin{theorem}[Divisor de puntos fijos]\label{thm:fixed-divisor}
Sea \(F\) una forma binaria de grado \(m\) sin factores repetidos, sobre un cuerpo de
característica que no divide \(m\). Entonces \(\cD_F\) es un morfismo de grado
\(m-1\) y sus puntos fijos son exactamente los ceros de \(F\).
\end{theorem}

\begin{proof}
Si \(F\) es libre de cuadrados, \(F_X\) y \(F_Y\) no tienen factor común. En efecto,
un factor lineal común dividiría, por Euler,
\[
XF_X+YF_Y=mF,
\]
y también dividiría ambas derivadas, lo que obligaría a que apareciera con
multiplicidad al menos dos en \(F\). Por tanto, el mapa tiene grado \(m-1\).

Un punto \([X:Y]\) es fijo si el vector \((-F_Y,F_X)\) es proporcional a
\((X,Y)\). Esto equivale a que su determinante sea cero:
\[
\det\begin{pmatrix}X&-F_Y\\Y&F_X\end{pmatrix}
=XF_X+YF_Y=mF(X,Y).
\]
Como \(m\ne0\) en el cuerpo, la condición es exactamente \(F(X,Y)=0\).
\end{proof}

\subsection{Multiplicadores iguales}

En dinámica compleja, si \(z_0\) es un punto fijo de un mapa racional \(R\), el
número \(R'(z_0)\) se llama multiplicador. Su módulo determina la estabilidad local:
atractor si es menor que uno, repulsor si es mayor que uno e indiferente si es igual
a uno.

\begin{theorem}[Multiplicador en cada raíz]\label{thm:root-multiplier}
Sea \(Q\) un polinomio de grado \(m\ge2\) con raíces simples. Para
\[
\Phi_Q(t)=t-m\frac{Q(t)}{Q'(t)},
\]
cada raíz \(\rho\) es un punto fijo con multiplicador
\begin{equation}\label{eq:multiplier}
\boxed{\Phi_Q'(\rho)=1-m.}
\end{equation}
En particular, para \(m=2\) el multiplicador es \(-1\), y para \(m\ge3\) todas las
raíces son repulsoras sobre \(\CC\).
\end{theorem}

\begin{proof}
Derivando,
\[
\Phi_Q'(t)
=1-m\frac{(Q')^2-QQ''}{(Q')^2}.
\]
En una raíz simple \(\rho\), \(Q(\rho)=0\) y \(Q'(\rho)\ne0\). Por tanto,
\[
\Phi_Q'(\rho)=1-m.
\]
\end{proof}

\subsection{Caracterización de Doyle--McMullen}

La siguiente unicidad explica por qué el mapa polar es inevitable una vez se fijan el
divisor de puntos fijos y los multiplicadores.

\begin{theorem}[Unicidad por puntos fijos y multiplicadores]\label{thm:uniqueness}
Sea \(Q\in\CC[t]\) mónico, de grado \(m\ge2\), con raíces simples y finitas. Existe
un único mapa racional de grado \(m-1\) que fija todas las raíces de \(Q\) y tiene
multiplicador \(1-m\) en cada una. Ese mapa es
\[
\Phi_Q(t)=t-m\frac{Q(t)}{Q'(t)}.
\]
\end{theorem}

\begin{proof}
Sea \(R=A/B\), con \(A,B\) coprimos y grados a lo sumo \(m-1\). Los puntos fijos
son los ceros de
\[
A(t)-tB(t).
\]
Este polinomio tiene grado a lo sumo \(m\) y se anula en las \(m\) raíces de \(Q\),
por lo que es proporcional a \(Q\). Multiplicando simultáneamente \(A\) y \(B\) por
una constante, podemos normalizar
\begin{equation}\label{eq:fixed-normalization}
A(t)-tB(t)=-mQ(t).
\end{equation}
Sea \(\rho\) una raíz. Como \(A(\rho)=\rho B(\rho)\),
\[
R'(\rho)-1
=\frac{A'(\rho)-\rho B'(\rho)-B(\rho)}{B(\rho)}.
\]
Al derivar \cref{eq:fixed-normalization}, el numerador es \(-mQ'(\rho)\). La
condición \(R'(\rho)=1-m\) implica
\[
-m=\frac{-mQ'(\rho)}{B(\rho)},
\]
y por tanto \(B(\rho)=Q'(\rho)\) para las \(m\) raíces. Como \(B\) y \(Q'\)
tienen grado a lo sumo \(m-1\), deben coincidir. Finalmente,
\[
A=tQ'-mQ.
\]
\end{proof}

\begin{researchbox}[Lectura dinámica]
El mapa de Newton ordinario
\[
N_Q(t)=t-\frac{Q(t)}{Q'(t)}
\]
hace superatractoras las raíces simples. El mapa polar
\[
\Phi_Q(t)=t-m\frac{Q(t)}{Q'(t)}
\]
hace que todas tengan el mismo multiplicador \(1-m\). No es un método numérico para
buscar raíces cuando \(m>2\); su función es codificar el divisor de raíces de forma
proyectivamente natural.
\end{researchbox}

\subsection{El caso cuadrático recupera la involución armónica}

\begin{theorem}[La involución armónica es el primer mapa polar]\label{thm:quadratic-polar-harmonic}
Sea \(Q(t)=(t-r)(t-s)\), con \(r\ne s\). Entonces
\begin{equation}\label{eq:Phi2-harmonic}
\Phi_Q(t)=t-2\frac{Q(t)}{Q'(t)}=h_{r,s}(t).
\end{equation}
Por tanto, la construcción parabólica de \cref{thm:parabola-harmonic} es el caso de
grado dos del mapa polar exterior.
\end{theorem}

\begin{proof}
Como
\[
Q(t)=t^2-(r+s)t+rs,
\qquad
Q'(t)=2t-(r+s),
\]
tenemos
\[
t-2\frac{Q(t)}{Q'(t)}
=\frac{t(2t-r-s)-2(t^2-(r+s)t+rs)}{2t-r-s}
=\frac{(r+s)t-2rs}{2t-r-s}.
\]
Multiplicando numerador y denominador por \(-1\), se obtiene
\[
\frac{2rs-(r+s)t}{r+s-2t}=h_{r,s}(t).
\]
\end{proof}

El hecho de que \(h_{r,s}\) sea una involución también se deduce de la unicidad:
una transformación de Möbius con dos puntos fijos y multiplicador \(-1\) en ellos
tiene orden dos.

% =============================================================================
\section{El hessiano como divisor crítico}
% =============================================================================

\subsection{La teoría clásica de invariantes entra en la dinámica}

El mapa polar está construido con las primeras derivadas de \(F\). Sus puntos
críticos están gobernados por las segundas derivadas, como en la teoría clásica del
hessiano de formas binarias \cite{Olver}.

\begin{definition}[Hessiano de una forma binaria]
Para una forma binaria \(F\), definimos
\begin{equation}\label{eq:hessian}
\Hess(F)=
\det\begin{pmatrix}
F_{XX}&F_{XY}\\
F_{XY}&F_{YY}
\end{pmatrix}
=F_{XX}F_{YY}-F_{XY}^2.
\end{equation}
Es una forma binaria de grado \(2m-4\) si \(\deg F=m\).
\end{definition}

\begin{theorem}[Divisor crítico hessiano]\label{thm:hessian-critical}
Sea \(F\) una forma binaria cuadrado-libre de grado \(m\ge2\). El divisor crítico del
mapa \(\cD_F=[-F_Y:F_X]\) es el divisor de ceros de \(\Hess(F)\), contado con
multiplicidad.
\end{theorem}

\begin{proof}
Un levantamiento homogéneo de \(\cD_F\) es
\[
\cH_F=(-F_Y,F_X).
\]
Su matriz jacobiana es
\[
J\cH_F=
\begin{pmatrix}
-F_{XY}&-F_{YY}\\
F_{XX}&F_{XY}
\end{pmatrix}.
\]
El determinante es
\[
\det(J\cH_F)=F_{XX}F_{YY}-F_{XY}^2=\Hess(F).
\]
La anulación de este jacobiano homogéneo describe precisamente la ramificación del
mapa proyectivizado. Su grado \(2m-4\) coincide con el número total de puntos críticos
de un mapa de grado \(m-1\), según Riemann--Hurwitz.
\end{proof}

\begin{proposition}[Ecuación afín de los puntos críticos]\label{prop:critical-affine}
Para \(F_Q(X,Y)=Y^mQ(X/Y)\),
\begin{equation}\label{eq:hessian-affine}
\Hess(F_Q)(t,1)
=(m-1)\left[mQ(t)Q''(t)-(m-1)Q'(t)^2\right].
\end{equation}
Además,
\begin{equation}\label{eq:Phi-derivative}
\Phi_Q'(t)
=\frac{mQ(t)Q''(t)-(m-1)Q'(t)^2}{Q'(t)^2}
\end{equation}
en la carta donde \(Q'(t)\ne0\).
\end{proposition}

\begin{proof}
Ya sabemos que
\[
F_X=Q',
\qquad
F_Y=mQ-tQ'.
\]
Derivando nuevamente,
\[
F_{XX}=Q'',
\qquad
F_{XY}=(m-1)Q'-tQ'',
\]
\[
F_{YY}=m(m-1)Q-2(m-1)tQ'+t^2Q''.
\]
Sustituyendo en \cref{eq:hessian} y cancelando los términos con \(t\), queda
\cref{eq:hessian-affine}. La fórmula para \(\Phi_Q'\) se obtiene por derivación
directa de \(t-mQ/Q'\).
\end{proof}

\begin{figure}[ht]
\centering
\begin{tikzpicture}[>=Latex,node distance=10mm and 12mm,every node/.style={align=center}]
\node[draw=Navy,fill=Ice,rounded corners,minimum width=38mm,minimum height=13mm] (F) {forma binaria $F$};
\node[draw=DeepBlue,fill=Sky,rounded corners,minimum width=43mm,minimum height=13mm,right=of F] (dF) {primeras derivadas\\$(-F_Y,F_X)$};
\node[draw=Gold,fill=Warm,rounded corners,minimum width=43mm,minimum height=13mm,right=of dF] (map) {mapa polar $\cD_F$};
\node[draw=Green,fill=Mint,rounded corners,minimum width=45mm,minimum height=13mm,below=of dF] (hess) {segundas derivadas\\$\Hess(F)$};
\node[draw=Brick,fill=white,rounded corners,minimum width=45mm,minimum height=13mm,right=of hess] (crit) {divisor crítico};
\draw[->,thick,DeepBlue] (F)--(dF);
\draw[->,thick,Gold] (dF)--(map);
\draw[->,thick,Green] (F)--(hess);
\draw[->,thick,Brick] (hess)--(crit);
\draw[->,thick,Slate,dashed] (map)--(crit);
\end{tikzpicture}
\caption{La forma binaria controla simultáneamente los puntos fijos mediante \(F=0\) y los puntos críticos mediante \(\Hess(F)=0\).}
\label{fig:hessian-dynamics}
\end{figure}

\subsection{Compatibilidad con el teorema de puntos fijos racionales}

Un mapa racional de grado \(d\) con puntos fijos simples y multiplicadores
\(\lambda_i\) satisface la identidad holomorfa \cite{Milnor}
\[
\sum_i\frac{1}{1-\lambda_i}=1.
\]
En nuestro caso, \(d=m-1\), hay \(m=d+1\) puntos fijos y todos tienen
\(\lambda_i=1-m=-d\). Por tanto,
\[
\sum_{i=1}^{m}\frac1{1-(1-m)}
=m\cdot\frac1m=1.
\]
La igualdad de multiplicadores no contradice la dinámica global; encaja exactamente
en la restricción universal sobre puntos fijos.

% =============================================================================
\section{Espacios de módulos y configuraciones de puntos}
% =============================================================================

\subsection{Una correspondencia de clases geométricas}

La covariancia de \cref{thm:covariance} y la unicidad de \cref{thm:uniqueness}
permiten describir un estrato especial del espacio de módulos de mapas racionales.

Sea \(\cC_m\) el conjunto de configuraciones no ordenadas de \(m\) puntos distintos
de \(\PP^1(\CC)\), consideradas módulo \(\PGL_2(\CC)\). Sea \(\cE_{m-1}\) el
conjunto de clases de conjugación de mapas racionales de grado \(m-1\) con \(m\)
puntos fijos distintos, todos de multiplicador \(1-m\).

\begin{theorem}[Correspondencia modular en puntos geométricos]\label{thm:moduli-bijection}
La asignación
\begin{equation}\label{eq:moduli-map}
\{F=0\}\longmapsto[\cD_F]
\end{equation}
induce una biyección natural
\begin{equation}\label{eq:moduli-bijection}
\boxed{\cC_m\simeq\cE_{m-1}.}
\end{equation}
En particular, para \(m\ge3\), ambas familias tienen dimensión \(m-3\).
\end{theorem}

\begin{proof}
La covariancia muestra que configuraciones proyectivamente equivalentes producen
mapas conjugados. El divisor de puntos fijos de \(\cD_F\) es exactamente el divisor
de ceros de \(F\), por \cref{thm:fixed-divisor}; por tanto, la asignación es
inyectiva sobre clases geométricas. Recíprocamente, un mapa perteneciente a
\(\cE_{m-1}\) tiene un divisor de \(m\) puntos fijos. La unicidad de
\cref{thm:uniqueness}, después de elegir una carta que evite \(\infty\), muestra que
el mapa es el polar asociado a ese divisor. La dimensión de configuraciones de
\(m\) puntos en \(\PP^1\), módulo el grupo tridimensional \(\PGL_2\), es \(m-3\).
\end{proof}

\begin{remark}
El teorema se formula deliberadamente como biyección de clases geométricas. Convertirla
en un isomorfismo preciso de espacios o pilas de módulos requiere controlar
estabilizadores, compactificaciones y configuraciones degeneradas; esa refinación
pertenece a una continuación del proyecto.
\end{remark}

\subsection{Grado tres: una única clase cuadrática}

Cuando \(m=3\), cualquier triple de puntos distintos de \(\PP^1\) es proyectivamente
equivalente a las raíces cúbicas de la unidad. Para
\[
F_0(X,Y)=X^3-Y^3,
\]
se tiene
\[
\cD_{F_0}([X:Y])=[3Y^2:3X^2]=[Y^2:X^2].
\]
En coordenada afín,
\[
z\longmapsto\frac1{z^2}.
\]

\begin{corollary}[Clasificación del caso cúbico]\label{cor:cubic-conjugacy}
Para todo polinomio cúbico con tres raíces distintas, el mapa
\[
\Phi_P(z)=z-3\frac{P(z)}{P'(z)}
\]
es conjugado por una transformación de Möbius a \(z\mapsto1/z^2\).
\end{corollary}

\begin{proof}
Una transformación de Möbius lleva las tres raíces a las raíces de \(X^3-Y^3\). La
covariancia de \cref{thm:covariance} transporta el mapa polar a
\(z\mapsto1/z^2\).
\end{proof}

\subsection{Grado cuatro: aparece el parámetro de razón doble}

Para cuatro puntos, la razón doble deja un parámetro modular. Normalicemos la
configuración como
\[
\{0,1,\lambda,\infty\},
\qquad \lambda\in\CC\setminus\{0,1\}.
\]
La forma binaria es
\[
F_\lambda(X,Y)=XY(X-Y)(X-\lambda Y).
\]
Un cálculo directo produce
\begin{align}
-F_{\lambda,Y}
&=-X^3+2(\lambda+1)X^2Y-3\lambda XY^2,\\
F_{\lambda,X}
&=3X^2Y-2(\lambda+1)XY^2+\lambda Y^3.
\end{align}
Por tanto, en la carta afín,
\begin{equation}\label{eq:quartic-normal-map}
\Phi_\lambda(z)=
\frac{-z^3+2(\lambda+1)z^2-3\lambda z}
{3z^2-2(\lambda+1)z+\lambda}.
\end{equation}
Fija \(0,1,\lambda,\infty\), todos con multiplicador \(-3\). El parámetro
\(\lambda\) se identifica sólo hasta las seis transformaciones anharmónicas que
corresponden a permutar cuatro puntos.

\begin{figure}[ht]
\centering
\begin{tikzpicture}[>=Latex]
\draw[Slate,thick] (-5,0)--(5,0);
\foreach \x/\lab/\col in {-3.7/{0}/Brick,-1.2/{1}/DeepBlue,1.8/{\lambda}/Gold,4.2/{\infty}/Green}{
  \fill[\col] (\x,0) circle (2.5pt);
  \node[above=2mm,font=\small\bfseries,text=\col] at (\x,0) {$\lab$};
}
\draw[decorate,decoration={brace,amplitude=7pt},Navy]
 (-3.7,-.35)--(1.8,-.35)
 node[midway,below=8pt,font=\small] {la razón doble conserva un parámetro};
\node[draw=Navy,fill=Ice,rounded corners,align=center] at (0,-1.65)
 {$\displaystyle \Phi_\lambda(z)=
\frac{-z^3+2(\lambda+1)z^2-3\lambda z}
{3z^2-2(\lambda+1)z+\lambda}$};
\end{tikzpicture}
\caption{A partir de cuatro raíces, la configuración ya no puede eliminarse completamente mediante una transformación proyectiva.}
\label{fig:quartic-moduli}
\end{figure}

% =============================================================================
\section{Raíces múltiples y versión ponderada}
% =============================================================================

\subsection{Cancelación del factor común}

Sea
\begin{equation}\label{eq:multiple-P}
P(t)=\prod_{i=1}^{s}(t-r_i)^{m_i},
\qquad
m_i\ge1,
\qquad
\sum_{i=1}^{s}m_i=n,
\end{equation}
con raíces distintas \(r_i\). Entonces
\[
\frac{P'(t)}{P(t)}=\sum_{i=1}^{s}\frac{m_i}{t-r_i}.
\]
El mapa afín formal sigue siendo
\begin{equation}\label{eq:weighted-Phi}
\Phi_P(t)=t-\frac{n}{\displaystyle\sum_{i=1}^{s}\frac{m_i}{t-r_i}}.
\end{equation}
Las derivadas homogéneas \(F_X,F_Y\) tienen el factor común
\[
\prod_{i=1}^{s}(X-r_iY)^{m_i-1}.
\]
Después de cancelarlo, se obtiene un mapa de grado \(s-1\), no de grado \(n-1\).

\begin{theorem}[Mapa polar de un divisor con multiplicidades]\label{thm:weighted-map}
Para \(P\) como en \cref{eq:multiple-P}, el mapa reducido asociado a
\(t-nP/P'\) tiene grado \(s-1\). Cada raíz distinta \(r_i\) es un punto fijo con
multiplicador
\begin{equation}\label{eq:weighted-multiplier}
\boxed{\Phi_P'(r_i)=1-\frac{n}{m_i}.}
\end{equation}
\end{theorem}

\begin{proof}
Escribamos \(G(t)=\prod_{i=1}^{s}(t-r_i)\). Entonces
\[
P'(t)=\frac{P(t)}{G(t)}
\sum_{i=1}^{s}m_i\prod_{j\ne i}(t-r_j).
\]
El factor \(P/G\) se cancela en numerador y denominador de \(t-nP/P'\). El
denominador restante tiene grado \(s-1\) y no se anula en ninguna \(r_i\), por lo
que no hay cancelaciones adicionales. Cerca de \(r_i\),
\[
\frac{P(t)}{P'(t)}=\frac{t-r_i}{m_i}+O((t-r_i)^2).
\]
Por tanto,
\[
\Phi_P(t)=r_i+\left(1-\frac{n}{m_i}\right)(t-r_i)+O((t-r_i)^2).
\]
\end{proof}

\begin{remark}
La multiplicidad deja de ser un dato invisible: se transforma en el multiplicador
dinámico. Una raíz con multiplicidad grande puede ser atractora; una raíz simple
tiene multiplicador \(1-n\). Si todas las raíces coinciden, \(s=1\) y el mapa se
reduce al mapa constante igual a esa raíz.
\end{remark}

\subsection{Apolaridad con multiplicidades}

Cuando las raíces se repiten, repetir el mismo hiperplano osculador no añade
condiciones lineales. La multiplicidad debe codificarse mediante jets, es decir,
mediante derivadas sucesivas de la parametrización.

Para un subespacio proyectivo \(W\subset\PP^n\), escribiremos
\[
W^\perp=\{[X]\in\PP^n:B_n(w,X)=0\text{ para todo }[w]\in W\}.
\]
La no degeneración de \(B_n\) implica que, si \(W\) tiene dimensión proyectiva
\(d\), entonces \(W^\perp\) tiene dimensión \(n-d-1\).

\begin{theorem}[Intersección apolar confluyente]\label{thm:confluent-apolar}
Sea
\[
P(t)=\prod_{i=1}^{s}(t-r_i)^{m_i},
\qquad r_i\ne r_j\ (i\ne j),
\qquad \sum_{i=1}^{s}m_i=n.
\]
Entonces
\begin{equation}\label{eq:confluent-intersection}
\boxed{
K_P=\bigcap_{i=1}^{s}
\left(\Osc_{r_i}^{m_i-1}(\nu_n)\right)^\perp.}
\end{equation}
Cada factor de la intersección tiene codimensión \(m_i\); las multiplicidades suman
\(n\), y la intersección resultante es el único punto proyectivo \(K_P\).
\end{theorem}

\begin{proof}
Un punto \(X\) pertenece a
\(\left(\Osc_{r_i}^{m_i-1}(\nu_n)\right)^\perp\) si y sólo si
\[
B_n(\nu_n^{(j)}(r_i),X)=0,
\qquad 0\le j\le m_i-1.
\]
Como la bilinearidad permite derivar dentro del primer argumento, estas condiciones
son equivalentes a que el polinomio
\[
F_X(t)=B_n(\nu_n(t),X)
\]
tenga un cero de orden al menos \(m_i\) en \(r_i\). Por tanto, si \(X\) pertenece
a todos los factores de \cref{eq:confluent-intersection}, entonces \(P\mid F_X\).
El polinomio \(F_X\) tiene grado a lo sumo \(n\) y no es idénticamente nulo para un
punto proyectivo \(X\); como \(P\) es mónico de grado \(n\), existe
\(\lambda\ne0\) tal que \(F_X=\lambda P\). La segunda identidad maestra da
\(P(t)=B_n(\nu_n(t),K_P)\), y la comparación de coeficientes produce
\(X=\lambda K_P\). Recíprocamente, la identidad maestra muestra que \(K_P\)
satisface todas las condiciones de orden de anulación.
\end{proof}

\begin{remark}
La formulación esquemática dice que un divisor de longitud \(n\) sobre la curva
racional normal determina, por apolaridad, un punto dual. El teorema anterior contiene
la versión lineal exacta que se necesita aquí; una compactificación modular exigiría
además organizar estas condiciones en familias cuando los puntos colisionan.
\end{remark}

% =============================================================================
\section{Una tentación geométrica falsa y su corrección}
% =============================================================================

Una investigación rigurosa no sólo debe demostrar lo correcto; también debe aislar
las generalizaciones seductoras que fallan.

En el caso \(o=0\), proyectar \(K_P\) desde \(\Osc_0^{n-2}\) hacia la tangente en
\(\infty\) produce \(-nP(0)/P'(0)\). Es tentador reemplazar \(0\) por un punto
variable \(o\) y concluir que la proyección desde \(\Osc_o^{n-2}\) produciría
\[
o-n\frac{P(o)}{P'(o)}.
\]
Eso no es correcto en las coordenadas tangentes usadas aquí.

\begin{proposition}[Resultado de la proyección móvil]\label{prop:false-generalization}
Con la identificación afín \([X_0:X_1]\mapsto X_0/X_1\) de la tangente
\(T_\infty\), la proyección de \(K_P\) desde \(\Osc_o^{n-2}(\nu_n)\) hacia
\(T_\infty\) produce
\begin{equation}\label{eq:moving-projection}
\Psi_P(o)=no-n\frac{P(o)}{P'(o)}
=(n-1)o+\Phi_P(o),
\end{equation}
no \(\Phi_P(o)\).
\end{proposition}

\begin{proofmap}[Idea de la verificación]
Trasladar \(o\) a \(0\) cambia el polinomio a \(P(o+u)\). La proyección en la
coordenada local \(u\) produce \(-nP(o)/P'(o)\). Sin embargo, la representación de
grado \(n\) actúa sobre la tangente en \(\infty\) con peso \(n\): una traslación
\(u\mapsto u+o\) suma \(no\), no \(o\), a la coordenada \(X_0/X_1\). Al deshacer la
traslación aparece \cref{eq:moving-projection}.
\end{proofmap}

Para \(n=2\), la proyección móvil produce \(o+h_{r,s}(o)\), mientras que el mapa
polar es \(h_{r,s}(o)\). Esta diferencia muestra por qué el mapa exterior y la familia
de proyecciones no deben identificarse sin controlar la representación tangente.

\begin{warning}
La coincidencia de dos construcciones en \(0\) y \(\infty\) no autoriza a suponer que
coinciden para un parámetro variable. La corrección anterior es importante: evita
convertir una observación verdadera en un teorema falso por extrapolación.
\end{warning}

% =============================================================================
\section{Ejemplos comparativos}
% =============================================================================

\subsection[Dos raíces: 1 y 4]{Dos raíces: \(1\) y \(4\)}

Para
\[
P(t)=(t-1)(t-4)=t^2-5t+4,
\]
se tiene
\[
a_0=1,
\qquad a_1=\frac52,
\qquad a_2=4.
\]
Por tanto,
\[
\cN_1=\frac52=\AM,
\qquad
\cN_2=\frac{4}{5/2}=\frac85=\HM,
\qquad
\cN_1\cN_2=4=\GM^2.
\]
El mapa polar es
\[
\Phi_P(t)=t-2\frac{t^2-5t+4}{2t-5}
=\frac{8-5t}{5-2t}=h_{1,4}(t).
\]
La parábola, la apolaridad y la dinámica racional describen exactamente la misma
involución.

\subsection{Cuatro raíces en progresión geométrica}

Tomemos \(r=(1,2,4,8)\). Entonces
\[
e_1=15,
\qquad e_2=70,
\qquad e_3=120,
\qquad e_4=64.
\]
Los coeficientes normalizados son
\[
a_0=1,
\quad a_1=\frac{15}{4},
\quad a_2=\frac{35}{3},
\quad a_3=30,
\quad a_4=64.
\]
La cadena queda
\[
\cN_1=\frac{15}{4}=3.75,
\qquad
\cN_2=\frac{28}{9}\approx3.111,
\qquad
\cN_3=\frac{18}{7}\approx2.571,
\qquad
\cN_4=\frac{32}{15}\approx2.133.
\]
Su producto es
\[
\frac{15}{4}\cdot\frac{28}{9}\cdot\frac{18}{7}\cdot\frac{32}{15}=64,
\]
y por tanto su media geométrica es \(64^{1/4}=2\sqrt2\), igual a la media geométrica
de \(1,2,4,8\).

\begin{figure}[ht]
\centering
\begin{tikzpicture}
\begin{axis}[
 ybar,width=0.88\textwidth,height=0.43\textwidth,
 ymin=0,ymax=4.25,bar width=18pt,
 symbolic x coords={N1,N2,N3,N4},
 xtick=data,xticklabels={$\cN_1$,$\cN_2$,$\cN_3$,$\cN_4$},ylabel={valor},
 nodes near coords,nodes near coords align={vertical},
 grid=major,grid style={SoftGray},
 every axis plot/.append style={fill=DeepBlue!45,draw=DeepBlue}]
\addplot coordinates {(N1,3.75) (N2,3.111111) (N3,2.571429) (N4,2.133333)};
\addplot[sharp plot,Brick,very thick,mark=none,nodes near coords={}] coordinates {(N1,2.828427) (N4,2.828427)};
\node[text=Brick,font=\small\bfseries] at (axis cs:N3,2.98) {$\GM=2\sqrt2$};
\end{axis}
\end{tikzpicture}
\caption{La cadena es decreciente. La media geométrica puede quedar entre peldaños; su papel exacto es multiplicativo, no necesariamente posicional.}
\label{fig:four-root-chain}
\end{figure}

\subsection{Una configuración simétrica compleja}

Sea
\[
P(t)=t^n-1.
\]
La forma binaria es \(F(X,Y)=X^n-Y^n\), y
\[
\cD_F([X:Y])=[Y^{n-1}:X^{n-1}].
\]
En coordenada afín,
\begin{equation}\label{eq:power-reciprocal}
\Phi_P(z)=\frac{1}{z^{n-1}}.
\end{equation}
Los puntos fijos son las raíces \(n\)-ésimas de la unidad, y cada multiplicador es
\(-(n-1)\). El hessiano es proporcional a
\[
X^{n-2}Y^{n-2},
\]
de modo que los únicos puntos críticos son \(0\) y \(\infty\), cada uno con
multiplicidad \(n-2\). Este ejemplo es el modelo de máxima simetría.

% =============================================================================
\section{Qué parte es clásica y qué parte constituye el programa de investigación}
% =============================================================================

\subsection{Inventario honesto}

\begin{longtable}{>{\RaggedRight\arraybackslash}p{0.28\textwidth}
                  >{\RaggedRight\arraybackslash}p{0.31\textwidth}
                  >{\RaggedRight\arraybackslash}p{0.32\textwidth}}
\toprule
\textbf{Componente} & \textbf{Estatus matemático} & \textbf{Papel en este manuscrito}\\
\midrule
\endfirsthead
\toprule
\textbf{Componente} & \textbf{Estatus matemático} & \textbf{Papel en este manuscrito}\\
\midrule
\endhead
Razón doble e involución armónica & Clásico en geometría proyectiva & Punto de entrada elemental y caso \(n=2\).\\
\addlinespace
Parábola que recupera \(\AM\) y \(\HM\) & Identidad elemental con lectura proyectiva & Realización visual rigurosa de \(h_{r,s}\).\\
\addlinespace
Curva racional normal y apolaridad & Teoría clásica de formas binarias & Produce el punto \(K_P\) y sus hiperplanos duales.\\
\addlinespace
Intersección de hiperplanos osculadores & Clásico; aparece en literatura moderna sobre curvas racionales normales & Se usa para interpretar \(K_P\) como compresión geométrica de raíces.\\
\addlinespace
Desigualdades de Newton--Maclaurin & Clásicas & Demuestran la monotonía de la cadena.\\
\addlinespace
Mapa \(t-nP/P'\) y unicidad por multiplicadores & Explícito en Doyle--McMullen y relacionado con el diferencial de formas binarias & Nivel superior de la torre polar.\\
\addlinespace
Divisor crítico hessiano & Clásico en teoría de invariantes y mapas gradiente & Conecta dinámica con el hessiano de \(F\).\\
\addlinespace
Cadena \(\cN_k=a_k/a_{k-1}\) como espectro de proyecciones de \(K_P\) & Síntesis específica desarrollada aquí; prioridad absoluta no afirmada & Une todas las medias entre \(\AM\) y \(\HM\).\\
\addlinespace
Identidad \(\cN_k=\Phi_{Q_k}(0)\) para la torre derivada & Deducción estructural central de este trabajo & Conecta apolaridad, derivadas y dinámica.\\
\addlinespace
Interpretación de \(\cN_k\) como \(\HM\) de raíces de \(P^{(n-k)}\) & Consecuencia exacta de la torre normalizada & Proporciona una lectura nueva y calculable de cada peldaño.\\
\addlinespace
Producto de las medias armónicas derivadas igual a \(\GM^n\) & Identidad telescópica dentro del nuevo marco & Integra la media geométrica sin forzarla a ser un peldaño.\\
\bottomrule
\end{longtable}

\subsection{Qué sería necesario para una publicación de investigación}

El manuscrito ya contiene teoremas completos y una arquitectura coherente. Para
convertirlo en un artículo de investigación sometible a una revista especializada,
serían necesarios al menos cuatro pasos adicionales:

\begin{enumerate}[leftmargin=9mm]
\item \textbf{Búsqueda bibliográfica de prioridad.} Rastrear de forma sistemática
MathSciNet, zbMATH y literatura clásica de invariantes para determinar si la cadena
\(a_k/a_{k-1}\), su interpretación mediante derivadas y su realización osculadora ya
aparecen conjuntamente.
\item \textbf{Formulación functorial.} Construir la torre \(Q_k\) y las proyecciones
\(\pi_k\) sin coordenadas, mediante fibrados de jets sobre \(\PP^1\), banderas de
osculación y cocientes de representaciones de \(\SL_2\).
\item \textbf{Teoría modular.} Elevar la biyección de \cref{thm:moduli-bijection} a
una afirmación sobre variedades o pilas, incluyendo divisores con multiplicidad y
compactificaciones estables.
\item \textbf{Dinámica crítica.} Estudiar cómo la configuración de raíces determina
las órbitas del divisor hessiano y clasificar las configuraciones para las cuales el
mapa polar es poscríticamente finito.
\end{enumerate}

% =============================================================================
\section{Direcciones de investigación y conjeturas prudentes}
% =============================================================================

\subsection{Problema A: caracterización axiomática de la cadena}

Las medias \(\cN_k\) poseen simetría, homogeneidad, idempotencia, monotonía,
dualidad recíproca y producto telescópico. Es natural preguntar si estas propiedades,
junto con una compatibilidad adecuada con la eliminación de variables o con la torre
derivada, caracterizan la cadena.

\begin{researchbox}[Problema A]
Clasificar las familias de medias
\(M_1,\dots,M_n\) sobre \((0,\infty)^n\) que satisfacen simultáneamente:
\[
M_1=\AM,
\qquad M_n=\HM,
\qquad M_1\ge\cdots\ge M_n,
\qquad \prod_{k=1}^{n}M_k=\GM^n,
\]
y una dualidad
\[
M_k(r^{-1})=M_{n-k+1}(r)^{-1}.
\]
¿Bajo qué axiomas adicionales resulta forzoso que \(M_k=\cN_k\)?
\end{researchbox}

\subsection{Problema B: una demostración puramente intercalatoria}

La monotonía \(\cN_k\ge\cN_{k+1}\) se demostró mediante Newton. La interpretación
como medias armónicas de raíces derivadas sugiere una prueba diferente.

\begin{researchbox}[Problema B]
Sea \(Q_{k+1}'=(k+1)Q_k\), con todos los ceros positivos y reales. Encontrar una
prueba directa, basada únicamente en intercalación o en transporte de medidas de
raíces, de
\[
\HM(\operatorname{Roots}Q_k)
\ge
\HM(\operatorname{Roots}Q_{k+1}).
\]
Una prueba de este tipo podría revelar un mecanismo geométrico detrás de las
desigualdades de Newton.
\end{researchbox}

\subsection{Problema C: dinámica del hessiano}

El divisor de raíces fija los puntos fijos; el hessiano fija los puntos críticos. La
dinámica global depende de cómo se relacionan ambos divisores.

\begin{researchbox}[Problema C]
Clasificar las formas binarias cuadrado-libres \(F\) para las cuales el mapa
\(\cD_F\) es poscríticamente finito. En particular:
\begin{enumerate}[leftmargin=8mm]
\item determinar cuándo los ceros de \(\Hess(F)\) tienen órbitas finitas;
\item describir el papel de los grupos finitos de simetría de la configuración de
raíces;
\item estudiar si existen familias no rígidas con relaciones críticas persistentes.
\end{enumerate}
\end{researchbox}

\subsection{Problema D: compactificación modular}

Cuando las raíces colisionan, el grado del mapa polar disminuye después de cancelar
factores comunes. Esto indica que la frontera natural del espacio modular no vive en
un único espacio \(\operatorname{Rat}_{n-1}\), sino en una unión estratificada por el
número de raíces distintas.

\begin{researchbox}[Problema D]
Construir una compactificación del locus \(\cE_{n-1}\) que registre simultáneamente:
\begin{itemize}[leftmargin=8mm]
\item colisiones de puntos en el divisor de raíces;
\item descenso de grado del mapa polar;
\item multiplicadores ponderados \(1-n/m_i\);
\item datos de osculación confluyente del punto \(K_P\).
\end{itemize}
\end{researchbox}

\subsection{Problema E: varias variables}

En dos variables, la forma de área identifica el diferencial de una forma binaria con
un vector y produce un auto-mapa de \(\PP^1\). En dimensión mayor, el gradiente de
una forma homogénea produce un mapa polar
\[
\PP^{d-1}\dashrightarrow(\PP^{d-1})^*,
\]
pero ya no existe una identificación canónica con el espacio original sin estructura
adicional.

\begin{researchbox}[Problema E]
Buscar análogos de la cadena \(\cN_k\) para polinomios hiperbólicos de varias
variables, usando gradientes, conos de hiperbicidad, polarización y desigualdades
lorentzianas. El objetivo no sería trasladar literalmente la media armónica, sino
identificar cocientes de coeficientes con significado geométrico y propiedades de
monotonía comparables.
\end{researchbox}

\begin{figure}[p]
\centering
\begin{tikzpicture}[>=Latex,node distance=11mm and 13mm,every node/.style={align=center}]
\node[draw=Navy,fill=Ice,rounded corners,minimum width=42mm,minimum height=14mm] (core)
 {núcleo probado\\cadena + torre polar};
\node[draw=Gold,fill=Warm,rounded corners,minimum width=43mm,minimum height=14mm,above left=of core] (axiom)
 {caracterización\\axiomática};
\node[draw=Green,fill=Mint,rounded corners,minimum width=43mm,minimum height=14mm,above right=of core] (inter)
 {intercalación\\y desigualdades};
\node[draw=Brick,fill=white,rounded corners,minimum width=43mm,minimum height=14mm,below left=of core] (dyn)
 {dinámica crítica\\del hessiano};
\node[draw=DeepBlue,fill=Sky,rounded corners,minimum width=43mm,minimum height=14mm,below right=of core] (mod)
 {módulos y\\compactificaciones};
\node[draw=Slate,fill=SoftGray,rounded corners,minimum width=43mm,minimum height=14mm,below=20mm of core] (multi)
 {formas hiperbólicas\\en varias variables};
\foreach \x in {axiom,inter,dyn,mod,multi}{\draw[-{Latex[length=2.5mm]},thick,Slate] (core)--(\x);}
\end{tikzpicture}
\caption{Agenda de investigación que emerge del núcleo demostrado. Las flechas señalan extensiones, no resultados ya establecidos.}
\label{fig:research-roadmap}
\end{figure}

% =============================================================================
\section{Conclusiones}
% =============================================================================

El recorrido comenzó con una construcción elemental en una parábola y terminó en una
red que conecta geometría proyectiva, polinomios simétricos, apolaridad, derivadas,
formas binarias y dinámica racional.

La primera conclusión es que la aparición simultánea de la media aritmética y la
media armónica no es accidental. En grado dos, ambas son imágenes de \(\infty\) y
\(0\) bajo la involución armónica fijada por las raíces. La parábola realiza esa
involución mediante secantes que pasan por el vértice.

La segunda conclusión es que el fenómeno no termina en dos raíces. El punto apolar
\[
K_P=[a_n:a_{n-1}:\cdots:a_0]
\]
comprime todos los polinomios simétricos normalizados. Las banderas osculadoras
opuestas en \(0\) y \(\infty\) aíslan cada ventana
\([a_k:a_{k-1}]\), y su coordenada es la media
\[
\cN_k=\frac{a_k}{a_{k-1}}.
\]
Las desigualdades de Newton convierten estas lecturas en una cadena decreciente desde
\(\AM\) hasta \(\HM\), mientras que el producto telescópico recupera \(\GM^n\).

La tercera conclusión es la existencia de una torre dinámica:
\[
Q_k=\frac{k!}{n!}P^{(n-k)},
\qquad
\Phi_k=t-\frac{Q_k}{Q_{k-1}}.
\]
Cada peldaño de la cadena satisface
\[
\cN_k=\Phi_k(0),
\]
y, para raíces positivas, es la media armónica de las raíces de \(Q_k\). Todas las
capas conservan la media aritmética original. Así, la cadena puede entenderse como
una secuencia de medias armónicas de configuraciones derivadas que se intercalan.

La cuarta conclusión es que el nivel superior
\[
\Phi_P=t-nP/P'
\]
es un objeto proyectivo clásico asociado al diferencial exterior de la forma binaria
de raíces. Sus puntos fijos son las raíces, todos con multiplicador \(1-n\), y sus
puntos críticos forman el divisor hessiano. Esta estructura produce una
correspondencia natural entre configuraciones de puntos en \(\PP^1\) y un locus de
mapas racionales con multiplicadores fijos.

Finalmente, el análisis mostró un principio metodológico: una identidad elegante no
debe confundirse con una novedad profunda, pero tampoco debe descartarse por ser
algebraicamente sencilla. Su valor puede estar en revelar que objetos clásicos,
considerados por separado, forman una arquitectura no evidente. La cadena apolar y la
torre polar constituyen precisamente esa arquitectura.

\Needspace{0.48\textheight}
\begin{keybox}[Teorema-síntesis]
Para raíces positivas \(r_1,\dots,r_n\), sea
\[
P(t)=\prod_{i=1}^{n}(t-r_i),
\qquad
Q_k(t)=\frac{k!}{n!}P^{(n-k)}(t),
\qquad
\Phi_k(t)=t-\frac{Q_k(t)}{Q_{k-1}(t)}.
\]
Entonces
\[
\boxed{
\AM=\Phi_1(0)\ge\Phi_2(0)\ge\cdots\ge\Phi_n(0)=\HM,}
\]
\[
\boxed{
\prod_{k=1}^{n}\Phi_k(0)=\GM^n,}
\]
y \(\Phi_k(0)\) es simultáneamente:
\begin{enumerate}[leftmargin=8mm]
\item el cociente \([a_k:a_{k-1}]\) leído desde el punto apolar \(K_P\);
\item la media armónica de las raíces de \(Q_k\);
\item la imagen del origen bajo el mapa polar exterior de la forma binaria asociada a
\(Q_k\).
\end{enumerate}
\end{keybox}

% =============================================================================
\FloatBarrier
\appendix
\section{Formulario de identidades}
% =============================================================================

\begin{align*}
P(t)&=\prod_{i=1}^{n}(t-r_i)
=\sum_{k=0}^{n}(-1)^ke_kt^{n-k},\\
a_k&=\frac{e_k}{\binom nk},\\
\cN_k&=\frac{a_k}{a_{k-1}}
=\frac{k}{n-k+1}\frac{e_k}{e_{k-1}},\\
\cN_1&=\AM,\\
\cN_n&=\HM,\\
\prod_{k=1}^{n}\cN_k&=\GM^n,\\
\cN_k(r^{-1})&=\cN_{n-k+1}(r)^{-1},\\
Q_k(t)&=\frac{k!}{n!}P^{(n-k)}(t),\\
Q_k'&=kQ_{k-1},\\
Q_k(0)&=(-1)^ka_k,\\
\Phi_k(t)&=t-k\frac{Q_k(t)}{Q_k'(t)}
=t-\frac{Q_k(t)}{Q_{k-1}(t)},\\
\Phi_k(0)&=\cN_k,\\
\Phi_P(t)&=t-n\frac{P(t)}{P'(t)},\\
\Phi_P'(r_i)&=1-n\qquad(r_i\text{ simple}),\\
\Phi_Q'(t)&=\frac{mQQ''-(m-1)(Q')^2}{(Q')^2},\\
\Hess(F_Q)(t,1)&=(m-1)\bigl[mQQ''-(m-1)(Q')^2\bigr].
\end{align*}

% =============================================================================
\section{Detalles computacionales verificables}
% =============================================================================

Las identidades principales son simbólicas y no dependen de experimentos. Sin
embargo, para verificar ejemplos o explorar conjeturas puede emplearse el siguiente
pseudocódigo algebraico.

\begin{verbatim}
Entrada: raíces positivas r[1],...,r[n]
1. Construir P(t) = producto_i (t-r[i]).
2. Para k=0,...,n:
       e[k] = coeficiente de (-1)^k t^(n-k) en P.
       a[k] = e[k] / binomial(n,k).
3. Para k=1,...,n:
       N[k] = a[k]/a[k-1].
       Q[k] = (k!/n!) * derivada(P,n-k).
       Phi[k](t) = t - Q[k](t)/Q[k-1](t).
4. Verificar:
       Phi[k](0) == N[k].
       N[k] >= N[k+1].
       producto_k N[k] == producto_i r[i].
5. Para cada k, hallar las raíces rho[k,*] de Q[k] y verificar:
       N[k] == media_armonica(rho[k,*]).
\end{verbatim}

% =============================================================================
\section{Referencias comentadas}
% =============================================================================

\begin{thebibliography}{99}

\bibitem{Berger}
M.~Berger,
\emph{Geometry I},
Springer, Berlín, 1987.
\newblock Referencia general para geometría clásica y proyectiva.

\bibitem{CaminataCarliniSchaffler}
A.~Caminata, E.~Carlini y L.~Schaffler,
\emph{Simplices Osculating Rational Normal Curves},
Vietnam Journal of Mathematics \textbf{54} (2026), 627--640; publicado en línea el 16 de abril de 2025.
\newblock Desarrolla configuraciones de hiperplanos osculadores de curvas racionales normales y ofrece un contexto moderno para la geometría usada aquí.

\bibitem{Coxeter}
H.~S.~M. Coxeter,
\emph{Projective Geometry}, segunda edición,
Springer, Nueva York, 1987.
\newblock Fuente clásica para razón doble, división armónica y transformaciones proyectivas.

\bibitem{Dolgachev}
I.~Dolgachev,
\emph{Classical Algebraic Geometry: A Modern View},
Cambridge University Press, Cambridge, 2012.
\newblock Tratamiento moderno de geometría algebraica clásica, formas binarias e invariantes.

\bibitem{DoyleMcMullen}
P.~Doyle y C.~McMullen,
\emph{Solving the Quintic by Iteration},
Acta Mathematica \textbf{163} (1989), 151--180. DOI: 10.1007/BF02392735.
\newblock En su discusión de mapas racionales con simetría aparece el mapa asociado al diferencial de una forma binaria y la caracterización de \(t-nP/P'\) por sus puntos fijos y multiplicadores.

\bibitem{GraceYoung}
J.~H. Grace y A.~Young,
\emph{The Algebra of Invariants},
Cambridge University Press, Cambridge, 1903; reimpresión Chelsea.
\newblock Fuente histórica de teoría clásica de invariantes y formas binarias.

\bibitem{HardyLittlewoodPolya}
G.~H. Hardy, J.~E. Littlewood y G.~Pólya,
\emph{Inequalities}, segunda edición,
Cambridge University Press, Cambridge, 1952.
\newblock Referencia clásica para desigualdades entre medias y polinomios simétricos.

\bibitem{Harris}
J.~Harris,
\emph{Algebraic Geometry: A First Course},
Springer, Nueva York, 1992.
\newblock Introducción a curvas racionales normales, dualidad y geometría proyectiva algebraica.

\bibitem{IarrobinoKanev}
A.~Iarrobino y V.~Kanev,
\emph{Power Sums, Gorenstein Algebras, and Determinantal Loci},
Lecture Notes in Mathematics 1721, Springer, Berlín, 1999.
\newblock Contexto profundo para apolaridad y álgebra de formas.

\bibitem{MarshallOlkinArnold}
A.~W. Marshall, I.~Olkin y B.~C. Arnold,
\emph{Inequalities: Theory of Majorization and Its Applications},
segunda edición, Springer, Nueva York, 2011.
\newblock Marco moderno para desigualdades simétricas y mayoración.

\bibitem{Milnor}
J.~Milnor,
\emph{Dynamics in One Complex Variable}, tercera edición,
Princeton University Press, Princeton, 2006.
\newblock Referencia estándar para dinámica de mapas racionales, puntos fijos, multiplicadores y puntos críticos.

\bibitem{Olver}
P.~J. Olver,
\emph{Classical Invariant Theory},
Cambridge University Press, Cambridge, 1999.
\newblock Presentación sistemática de invariantes, covariantes, hessianos y formas binarias.

\bibitem{Silverman}
J.~H. Silverman,
\emph{The Arithmetic of Dynamical Systems},
Springer, Nueva York, 2007.
\newblock Contexto aritmético y modular de los mapas racionales.

\end{thebibliography}

\vfill
\begin{center}
\begin{tikzpicture}
\draw[Gold,line width=0.9pt] (0,0)--(10,0);
\node[below=3mm,font=\small\sffamily,text=Slate,align=center] at (5,0)
 {Fin del manuscrito\\Versión de investigación, julio de 2026};
\end{tikzpicture}
\end{center}

\end{document}
